%% =====================================================================
%%  Capítulo 9 -- Sistemas variantes en el tiempo
%%  Formato nuevo (noviembre de 2026). Integra: diapositivas
%%  Present2026/Clase101926_VariantesTiempo (y su versión StmasVariantes_Inestab2026,
%%  exportada como VariantesTp_Inestabilidad.pdf), Ecuaciones_LaTeXiT.tex, el
%%  "Solucionario y guion" (19 de octubre de 2026), VanDerPol_Inestabilidad.pptx y los
%%  scripts Codigos/MATLAB_VariantesTiempo/TV01--TV09, reescritos en codigo/cap09_variantes.
%%  Handout6a.tex era de preliminares matemáticos (ya en el capítulo 3).
%%  Barbalat y LaSalle--Yoshizawa se prueban aquí (§9.6); los conversos están en el capítulo 10.
%%  Correcciones: ver MIGRACION.md. Números verificados con
%%  herramientas/cap09_numeros.py y codigo/cap09_variantes/verificarCap09.m.
%% =====================================================================

\capitulo[cap:variantes]{Sistemas variantes en el tiempo}
  {Nunca, ni en las situaciones más desesperadas, he perdido mi sentido del humor\ldots}
  {Enrique Vila-Matas}

\begin{objetivos}[partes/09-variantes.html]
  \item Explicar por qué la solución de $\dot x=f(t,x)$ depende del instante inicial y distinguir estabilidad, estabilidad uniforme, estabilidad asintótica y estabilidad asintótica uniforme.
  \item Usar las funciones de comparación ($\K$, $\Kinf$, $\KL$) para escribir cotas uniformes, y reconocer cuándo una función no es de esas clases.
  \item Aplicar los teoremas de Lyapunov con $V(t,x)$: verificar que $V$ es definida positiva, decreciente y radialmente no acotada, y calcular $\Vdot$ sin olvidar $\partial V/\partial t$.
  \item Reconocer los errores clásicos de los sistemas variantes: valores propios ``congelados'', linealización no uniforme y $\Vdot<0$ sin una cota independiente de $t$.
  \item Cerrar la brecha de $\Vdot\le0$ en sistemas variantes (LaSalle periódico, Floquet, LaSalle--Yoshizawa con Barbalat) y probar inestabilidad con los teoremas de inestabilidad y de Chetaev.
\end{objetivos}

%% ---------------------------------------------------------------------
\section{Por qué sistemas variantes en el tiempo}
\label{sec:vt-porque}
%% ---------------------------------------------------------------------

Hasta aquí casi todo el curso trató sistemas invariantes en el tiempo, $\dot x=f(x)$. En control, sin embargo, casi todo problema no lineal termina siendo variante en el tiempo, por tres caminos:
\begin{itemize}
  \item \textbf{Seguimiento.} Aunque la planta $\dot x=f(x,u)$ sea invariante, el error respecto a una referencia que se mueve, $e=x-x_{\mathrm{ref}}(t)$, obedece
  \begin{equation}
    \dot e=f\bigl(e+x_{\mathrm{ref}}(t),u\bigr)-\dot x_{\mathrm{ref}}(t)=:\bar f(t,e) ,
    \label{eq:vt-seguimiento}
  \end{equation}
  que depende explícitamente de $t$. Estabilizar el seguimiento es estabilizar el origen de un sistema variante.
  \item \textbf{Excitación paramétrica.} Un parámetro que oscila (la rigidez de un resorte, la longitud de un columpio, una ganancia) produce ecuaciones como la de Mathieu amortiguada, $\ddot y+\dot y+(2+\sin t)\,y=0$ (\cref{ej:vt-mathieu}).
  \item \textbf{Adaptación.} En control adaptativo, los parámetros estimados y las señales de excitación entran como funciones del tiempo: $\dot e=-e+\theta\,\omega(t)$, $\dot\theta=-e\,\omega(t)$ (\cref{ej:av-adaptativo}).
\end{itemize}

\begin{pausa}
Antes de seguir, prediga: (1) si $A(t)$ tiene todos sus valores propios con parte real negativa para cada $t$, ¿es estable $\dot x=A(t)x$? (2) Si $V(t,x)>0$ y $\Vdot(t,x)<0$ para todo $x\neq0$, ¿es asintóticamente estable el origen? (3) Si todas las soluciones tienden a cero, ¿tienden a cero ``igual de rápido'' sin importar cuándo empiezan?
\respuesta{Las tres respuestas son \emph{no}, y las tres tienen la misma estructura: algo que vale para sistemas invariantes en el tiempo y que se extrapola sin pensar. (1) Markus--Yamabe: valores propios $-\tfrac14\pm j\tfrac{\sqrt7}4$ para todo $t$ y soluciones que crecen como $e^{t/2}$ (\cref{ej:vt-markus}). (2) Con $\dot x=-x/(1+t)^2$ y $V=x^2$, $\Vdot<0$ fuera del origen y $x(t)\to x_0e^{-1/(1+t_0)}\neq0$ (\cref{ej:vt-e1}). (3) Con $\dot x=-x/(1+t)$, el tiempo para bajar de $0.1$ es $9(1+t_0)$: $9$ si se arranca en $t_0=0$, $909$ si se arranca en $t_0=100$ (\cref{ej:vt-motivador}).}
\end{pausa}

\begin{figure}[ht]
\centering
\includegraphics[width=0.95\linewidth]{cap09/t_mapa_variantes}
\caption{Mapa de las nociones de estabilidad para $\dot x=f(t,x)$ (\cref{sec:vt-uniformidad,sec:vt-comparacion}). Cada flecha es una implicación; las que no aparecen son falsas en general.}
\label{fig:vt-mapa}
\end{figure}

\notamargen{\textbf{Ruta de las clases.} Semana~11: uniformidad y funciones de comparación (\S\ref{sec:vt-uniformidad}--\ref{sec:vt-comparacion}) y teoremas de Lyapunov con $V(t,x)$ (\S\ref{sec:vt-lyapunov}). Semana~12: ejemplos, $\Vdot\le0$ con Barbalat y LaSalle--Yoshizawa, e inestabilidad (\S\ref{sec:vt-ejemplos}--\ref{sec:vt-inestabilidad}).}

\notamargen{\textbf{Vocabulario.} \citet{Kha02} y \citet{KeBr23} llaman \emph{autónomo} al sistema $\dot x=f(x)$ y \emph{no autónomo} a $\dot x=f(t,x)$. En el curso: invariante y variante en el tiempo.}

\begin{historia}[De Lyapunov a Kamke, Hahn y Yoshizawa]
Lyapunov trató desde el principio el caso variante en el tiempo: su tesis de 1892~\citep{Lya92} estudia $\dot x=f(t,x)$ y ya distingue funciones $V(t,x)$ que admiten una cota superior. Las nociones uniformes se precisaron en las décadas siguientes, y Massera dio en 1949 el primer teorema recíproco, para sistemas invariantes y periódicos~\citep{Mas49}; en 1956 lo extendió a los sistemas variantes en el tiempo, donde la propiedad que tiene recíproco es la estabilidad asintótica uniforme~\citep{Mas56}. Las funciones de comparación se llaman también funciones de Kamke, por Erich Kamke y sus teoremas de comparación para ecuaciones diferenciales; su uso sistemático en estabilidad se consolida con la monografía de W.~Hahn~\citep{Hah67}. Yoshizawa extendió en 1966 el argumento de LaSalle al caso variante~\citep{Yos66}. Markus y Yamabe publicaron en 1960 el ejemplo de los valores propios ``congelados''~\citep{MaYa60}, y N.~G.~Chetaev formuló en los años treinta el teorema de inestabilidad que lleva su nombre, recogido en su libro~\citep{Che61}.
\end{historia}

%% ---------------------------------------------------------------------
\section{Uniformidad: cuando la hora de arranque importa}
\label{sec:vt-uniformidad}
%% ---------------------------------------------------------------------

En todo el capítulo se considera
\begin{equation}
  \dot x=f(t,x),\qquad f:[0,\infty)\times\D\to\R^n ,
  \label{eq:vt-sistema}
\end{equation}
con $\D\subset\R^n$ un dominio que contiene al origen, $f$ continua a trozos en $t$ y localmente Lipschitz en $x$ uniformemente en $t$ (\cref{cap:preliminares}), y $f(t,0)=0$ para todo $t\ge0$: el origen es un equilibrio en todo instante. Se escribe $\phi(t;t_0,x_0)$ para la solución que en $t=t_0$ pasa por $x_0$. La diferencia con el caso invariante es justamente ese $t_0$: en $\dot x=f(x)$, $\phi(t;t_0,x_0)=\flujo{t-t_0}{x_0}$ y el instante inicial solo traslada la solución; aquí la cambia.

\notamargen{¿Por qué un equilibrio en el origen? Con el cambio $z=x-x^*(t)$ un equilibrio, o una solución nominal cualquiera $x^*(t)$, se traslada al origen; el precio es que el sistema en $z$ es variante en el tiempo aun si el original no lo era, como en \eqref{eq:vt-seguimiento}.}

\begin{ejemplo}[label=ej:vt-motivador]{El mismo sistema, dos horas de arranque}
Considere $\dot x=-\dfrac{x}{1+t}$. Calcule el tiempo que tarda $x$ en pasar de $x_0=1$ a $0.1$ si se arranca en $t_0=0$ y si se arranca en $t_0=100$.
\solucion Separando variables, $\ln\dfrac{x(t)}{x_0}=-\ln\dfrac{1+t}{1+t_0}$, es decir
\[
  x(t)=x_0\,\frac{1+t_0}{1+t} .
\]
$x(t)=0.1\,x_0$ cuando $1+t=10(1+t_0)$, así que el tiempo de asentamiento es
\[
  T(0.1,t_0)=t-t_0=9\,(1+t_0):\qquad T=9\ \text{si } t_0=0,\qquad T=909\ \text{si } t_0=100 .
\]
Todas las soluciones tienden a cero y $\abs{x(t)}\le\abs{x_0}$ para todo $t\ge t_0$, cualquiera que sea $t_0$. Pero la tasa $1/(1+t)$ se debilita con el reloj: cuanto más tarde se arranca, más lento se converge, y no hay un tiempo $T$ que sirva para todos los $t_0$ (\cref{fig:vt-tiempo}).
\end{ejemplo}

\begin{figure}[ht]
\centering
\includegraphics[width=\linewidth]{cap09/f_tiempo_inicial}
\caption{$\dot x=-x/(1+t)$ con $x_0=1$. Izquierda: soluciones para varios $t_0$; la línea punteada es el umbral $0.1$, que la solución con $t_0=0$ cruza en $t=9$ y la de $t_0=40$ cruza en $t=409$. Derecha: las mismas soluciones contra $s=t-t_0$. Cada una es una cota $\beta_{t_0}(1,s)=(1+t_0)/(1+t_0+s)$ de clase $\KL$, pero la envolvente de todas es la recta $1$: no hay una cota común que decaiga.}
\label{fig:vt-tiempo}
\end{figure}

\codigorepo{cap09_variantes/C9_Ej1_TiempoInicial.m}{tiempoinicial}{El tiempo de asentamiento $9(1+t_0)$, simulado}

\video{V9_1_tiempo_inicial}{Misma condición inicial, distinta hora de arranque}

\subsection{Las cuatro definiciones}

\begin{definicion}[label=def:vt-estabilidad]{Estabilidad de $\dot x=f(t,x)$ \citep[Def.~4.4]{Kha02}}
El origen de \eqref{eq:vt-sistema} es
\begin{itemize}
  \item \clave{estable} si para cada $\varepsilon>0$ y cada $t_0\ge0$ existe $\delta=\delta(\varepsilon,t_0)>0$ tal que
  \[
    \norm{x(t_0)}<\delta\implica\norm{x(t)}<\varepsilon\quad\text{para todo }t\ge t_0 ;
  \]
  \item \clave{uniformemente estable} (US) si, para cada $\varepsilon>0$, el $\delta$ anterior se puede tomar \emph{independiente de $t_0$}, $\delta=\delta(\varepsilon)$;
  \item \clave{inestable} si no es estable;
  \item \clave{asintóticamente estable} (AS) si es estable y existe $c=c(t_0)>0$ tal que $x(t)\to0$ cuando $t\to\infty$ para todo $\norm{x(t_0)}<c$;
  \item \clave{uniformemente asintóticamente estable} (UAS) si es US y existe $c>0$, \emph{independiente de $t_0$}, tal que para cada $\eta>0$ hay un $T=T(\eta)>0$ con
  \[
    \norm{x(t_0)}<c\implica\norm{x(t)}<\eta\quad\text{para todo }t\ge t_0+T(\eta)\ \text{y todo }t_0\ge0 ;
  \]
  \item \clave{globalmente uniformemente asintóticamente estable} (GUAS) si es US, $\delta(\varepsilon)$ se puede escoger con $\delta(\varepsilon)\to\infty$ cuando $\varepsilon\to\infty$, y para cada par $\eta,c>0$ hay un $T=T(\eta,c)$ con $\norm{x(t)}<\eta$ para todo $t\ge t_0+T$, todo $t_0\ge0$ y todo $\norm{x(t_0)}<c$.
\end{itemize}
\end{definicion}

Los dos requisitos de uniformidad son de naturaleza distinta. US dice que el \emph{tamaño} de la excursión no depende de la hora de arranque. UAS dice además que la \emph{velocidad} de la convergencia (el tiempo $T(\eta)$ para entrar a la bola $\Bola\eta$) tampoco depende de ella. En el \cref{ej:vt-motivador} se cumple lo primero con $\delta=\varepsilon$, y falla lo segundo: $T(0.1,t_0)=9(1+t_0)$ no tiene cota uniforme. El origen es US y AS, pero no UAS. Es el ejemplo de la parte baja del mapa de la \cref{fig:vt-mapa}.

\begin{trampa}[label=tr:vt-converge]{``Si toda solución converge a cero, la convergencia es uniforme''}
\plausible{En un sistema invariante en el tiempo, estable y atractivo es lo mismo que uniformemente asintóticamente estable: el instante inicial no cambia nada.}
En el \cref{ej:vt-motivador} toda solución converge y el origen es incluso uniformemente estable, pero el tiempo para entrar a $\Bola{0.1}$ crece sin cota con $t_0$. Tampoco es un detalle técnico: un controlador cuyo error converge ``a la larga'' pero cada vez más despacio no es aceptable, y las cotas uniformes son las que sobreviven a perturbaciones (los teoremas recíprocos de la \cref{sec:av-conversos} existen para UAS, no para AS).
\correcto{Converger y converger uniformemente son dos propiedades distintas. Para un sistema variante en el tiempo hay que verificar que $\delta$ y $T$ no dependan de $t_0$, o, equivalentemente, encontrar una sola cota $\beta\in\KL$ (\cref{lem:vt-kl}).}
\end{trampa}

\begin{pausa}
Para $\dot x=-x/(1+t)$, ¿qué pasa si se cambia la tasa por $-x/(1+t)^{1/2}$? ¿Y por $-(2+\sin t)\,x$?
\respuesta{Con $-x/\sqrt{1+t}$, $x(t)=x_0\exp\bigl(-2(\sqrt{1+t}-\sqrt{1+t_0})\bigr)$: la convergencia es más rápida que con $1/(1+t)$, pero $T(0.1,t_0)$ sigue creciendo con $t_0$ (como $\sqrt{t_0}$); de nuevo AS sin UAS. Con $-(2+\sin t)x$ la tasa está entre $1$ y $3$ para todo $t$, así que $\abs{x(t)}\le\abs{x_0}e^{-(t-t_0)}$: UAS, de hecho exponencial y global. Lo que destruye la uniformidad no es que la tasa varíe, sino que se \emph{apague}.}
\end{pausa}

\begin{trampa}[label=tr:vt-estable]{``Estable con $\delta(\varepsilon,t_0)$ es lo mismo que uniformemente estable''}
\plausible{Se dio un $\delta$ para cada $t_0$; tomar el menor de ellos debería servir para todos.}
El ínfimo de infinitos números positivos puede ser cero. Para $\dot x=(6t\sin t-2t)x$ \citep[\S4.5]{Kha02}, $x(t)=x(t_0)\exp\bigl(\int_{t_0}^t(6\tau\sin\tau-2\tau)\dif\tau\bigr)$ y el exponente, que es negativo a la larga, alcanza máximos que crecen con $t_0$ cuando $t_0=2n\pi$: el origen es AS, pero el $\delta$ necesario se encoge sin cota y no es US.
\correcto{US exige exhibir un $\delta(\varepsilon)$ que no dependa de $t_0$. En la práctica se prueba con una función de Lyapunov \emph{decreciente} (\cref{def:vt-decreciente}).}
\end{trampa}

En un sistema invariante en el tiempo todas las propiedades anteriores son automáticamente uniformes: si el origen de $\dot x=f(x)$ es estable, es US, y si es AS, es UAS. Por eso los capítulos anteriores no las distinguían. Lo mismo vale para un sistema \emph{periódico} en $t$, $f(t+\tau,x)=f(t,x)$: basta revisar $t_0\in[0,\tau]$, que es un compacto \citep[\S4.5]{Kha02}.

%% ---------------------------------------------------------------------
\section{Funciones de comparación}
\label{sec:vt-comparacion}
%% ---------------------------------------------------------------------

Las definiciones con $\varepsilon$, $\delta$ y $T$ son difíciles de manipular. La herramienta que las compacta, y que se usa en el resto del curso (\cref{cap:avanzada}), son las funciones de comparación: funciones escalares que miden el tamaño de la cota en función del tamaño de la condición inicial y del tiempo transcurrido.

\begin{definicion}[label=def:vt-clases]{Clases $\K$, $\Kinf$ y $\KL$ \citep[Def.~4.2 y 4.3]{Kha02}}
\begin{itemize}
  \item $\alpha:[0,a)\to[0,\infty)$ es de \clave{clase $\K$} si es continua, \emph{estrictamente} creciente y $\alpha(0)=0$.
  \item Es de \clave{clase $\Kinf$} si además $a=\infty$ y $\alpha(r)\to\infty$ cuando $r\to\infty$.
  \item $\beta:[0,a)\times[0,\infty)\to[0,\infty)$ es de \clave{clase $\KL$} si para cada $s$ fijo, $r\mapsto\beta(r,s)$ es de clase $\K$, y para cada $r$ fijo, $s\mapsto\beta(r,s)$ es decreciente y $\beta(r,s)\to0$ cuando $s\to\infty$.
\end{itemize}
\end{definicion}

Ejemplos: $\tan^{-1}r\in\K$ pero no $\Kinf$ (se satura en $\pi/2$); $r^c\in\Kinf$ para $c>0$; $\min\{r,1\}$ no es de clase $\K$ (es constante para $r\ge1$); $re^{-s}$, $r^c e^{-s}$ y $\dfrac{r}{sr+1}$ son de clase $\KL$. La primera variable es el tamaño de la condición inicial; la segunda, el tiempo transcurrido.

\begin{pausa}
¿Cuáles son de clase $\KL$? (a) $\dfrac{r^2}{1+s}$; (b) $re^{-sr}$; (c) $\min\{r,1\}\,e^{-s}$; (d) $\dfrac r{1+s}+re^{-s}$.
\respuesta{(a) Sí: $\partial\beta/\partial r=2r/(1+s)>0$ y en $s$ decrece a cero. (b) No: para $s>0$, $\partial\beta/\partial r=e^{-sr}(1-sr)<0$ si $r>1/s$; la función sube hasta $r=1/s$ y luego baja. (c) No: es constante en $r$ para $r\ge1$, y ``creciente'' no es ``estrictamente creciente''. (d) Sí: es suma de dos funciones de clase $\KL$.}
\end{pausa}

\begin{trampa}[label=tr:vt-kl]{``$re^{-sr}$ decae a cero, luego es de clase $\KL$''}
\plausible{Para cada $r$, $re^{-sr}\to0$ cuando $s\to\infty$, que es lo que se le pide a una cota de convergencia.}
El requisito que falla es el de la primera variable: con $s=2$, el máximo está en $r=1/2$, y una condición inicial más grande recibiría una cota \emph{más pequeña} (\cref{fig:vt-comparacion}, derecha). Una cota así no garantiza nada sobre la estabilidad, porque no se hace pequeña cuando $r$ se hace pequeño de manera uniforme en $s$.
\correcto{Hay que revisar las dos variables: estrictamente creciente en $r$ para cada $s$, y decreciente a cero en $s$ para cada $r$. Para (b) hay arreglo: como $e^{-y}\le1/(1+y)$ para $y\ge0$, se tiene $re^{-sr}\le r/(1+sr)$, que sí es de clase $\KL$. En estabilidad basta con que exista \emph{alguna} envolvente de clase $\KL$.}
\end{trampa}

\begin{figure}[ht]
\centering
\includegraphics[width=\linewidth]{cap09/f_comparacion}
\caption{Izquierda: $\tan^{-1}r\in\K$ (se satura), $r^2/10\in\Kinf$ y $0.9\min\{r,1\}\notin\K$. Centro: cortes de $r/(sr+1)\in\KL$ para varios $s$: cada corte crece con $r$ y baja con $s$. Derecha: cortes de $re^{-sr}$; con $s>0$ cada uno tiene un máximo en $r=1/s$ y no es de clase $\K$.}
\label{fig:vt-comparacion}
\end{figure}

\codigorepo{cap09_variantes/C9_Ej2_FuncionesComparacion.m}{comparacion}{Revisar numéricamente las dos condiciones de la clase $\KL$}

\subsection{Propiedades que se usan}

\begin{lema}[label=lem:vt-propiedades]{Álgebra de las funciones de comparación \citep[Lema~4.2]{Kha02}}
Sean $\alpha_1,\alpha_2\in\K$ en $[0,a)$, $\alpha_3,\alpha_4\in\Kinf$ y $\beta\in\KL$. Entonces:
\begin{enumerate}
  \item $\alpha_1^{-1}$ existe en $[0,\alpha_1(a))$ y es de clase $\K$; $\alpha_3^{-1}\in\Kinf$.
  \item $\alpha_1\circ\alpha_2\in\K$ y $\alpha_3\circ\alpha_4\in\Kinf$.
  \item $\sigma(r,s)=\alpha_1\bigl(\beta(\alpha_2(r),s)\bigr)$ es de clase $\KL$ (con dominios compatibles).
\end{enumerate}
\end{lema}

\begin{demostracion}
(1) y (2) se siguen de que la inversa y la composición de funciones continuas estrictamente crecientes que se anulan en cero lo son. Para (3), $\sigma$ es continua y $\sigma(0,s)=\alpha_1(\beta(0,s))=\alpha_1(0)=0$. Con $s$ fijo y $r_1<r_2$, la cadena $\alpha_2(r_1)<\alpha_2(r_2)\Rightarrow\beta(\alpha_2(r_1),s)<\beta(\alpha_2(r_2),s)\Rightarrow\sigma(r_1,s)<\sigma(r_2,s)$ usa que las tres funciones son \emph{estrictamente} crecientes en su primer argumento. Con $r$ fijo, $\beta(\alpha_2(r),\cdot)$ decrece a cero y $\alpha_1$, continua y nula en cero, conserva el orden y el límite.
\end{demostracion}

El paso (3) se lee de adentro hacia afuera: tamaño inicial $r$, se agranda con $\alpha_2$, se deja decaer con $\beta$ y se traduce con $\alpha_1$. Es exactamente la cota que aparece en la prueba del \cref{thm:vt-lyapunov}, y es la razón por la que se exige crecimiento \emph{estricto}: hay que invertir $\alpha_1$.

\begin{lema}[label=lem:vt-sandwich]{Funciones definidas positivas entre dos funciones $\K$ \citep[Lema~4.3]{Kha02}}
Si $W:\Bola r\to\R$ es continua y definida positiva, existen $\alpha_1,\alpha_2\in\K$ en $[0,r)$ con
\[
  \alpha_1(\norm x)\le W(x)\le\alpha_2(\norm x)\qquad\text{para todo }x\in\Bola r .
\]
Si $W$ está definida en $\R^n$ y es radialmente no acotada, se pueden tomar $\alpha_1,\alpha_2\in\Kinf$.
\end{lema}

La idea de la prueba es construir $\psi(s)=\min_{s\le\norm x\le r}W(x)$, que es continua, positiva y creciente, y quedarse por debajo con una función estrictamente creciente; por arriba, igual con $\max_{\norm x\le s}W(x)$. Para una forma cuadrática no hay que construir nada: $\lmin(P)\norm x^2\le x\T Px\le\lmax(P)\norm x^2$.

\begin{lema}[label=lem:vt-flujo]{La ecuación de comparación genera una función $\KL$ \citep[Lema~4.4]{Kha02}}
Sea $\alpha\in\K$ localmente Lipschitz en $[0,a)$. La ecuación escalar
\[
  \dot y=-\alpha(y),\qquad y(t_0)=y_0\in[0,a)
\]
tiene solución única para todo $t\ge t_0$, y $y(t)=\sigma(y_0,t-t_0)$ con $\sigma\in\KL$ en $[0,a)\times[0,\infty)$.
\end{lema}

Tres casos que se usan todo el tiempo:
\[
  \dot y=-ky\ \Rightarrow\ \sigma=re^{-ks},\qquad
  \dot y=-ky^2\ \Rightarrow\ \sigma=\frac r{krs+1},\qquad
  \dot y=-y^3\ \Rightarrow\ \sigma=\frac r{\sqrt{1+2r^2s}} .
\]
Toda función generada así cumple $\sigma(r,0)=r$ y $\sigma(\sigma(r,s_1),s_2)=\sigma(r,s_1+s_2)$: es un \emph{flujo}. No toda función $\KL$ lo es (\cref{prob:vt-flujo}). Junto con el lema de comparación (\cref{lem:pm-comparacion}), esto da la herramienta principal: si $\dot v\le-\alpha(v)$ y $v(t_0)\le y_0$, entonces $v(t)\le\sigma(y_0,t-t_0)$.

\subsection{La caracterización por funciones de comparación}

\begin{lema}[label=lem:vt-kl]{Estabilidad uniforme con funciones de comparación \citep[Lema~4.5]{Kha02}}
El origen de \eqref{eq:vt-sistema} es
\begin{itemize}
  \item US si y solo si existen $\alpha\in\K$ y $c>0$, independientes de $t_0$, tales que
  \begin{equation}
    \norm{x(t)}\le\alpha\bigl(\norm{x(t_0)}\bigr)\qquad\text{para todo }t\ge t_0\ge0\text{ y }\norm{x(t_0)}<c ;
    \label{eq:vt-us}
  \end{equation}
  \item UAS si y solo si existen $\beta\in\KL$ y $c>0$, independientes de $t_0$, tales que
  \begin{equation}
    \norm{x(t)}\le\beta\bigl(\norm{x(t_0)},\,t-t_0\bigr)\qquad\text{para todo }t\ge t_0\ge0\text{ y }\norm{x(t_0)}<c ;
    \label{eq:vt-uas}
  \end{equation}
  \item GUAS si y solo si \eqref{eq:vt-uas} vale para toda $x(t_0)\in\R^n$.
\end{itemize}
\end{lema}

\begin{ideaclave}
Una sola función $\beta\in\KL$ para todos los instantes iniciales: eso es la uniformidad. La primera variable de $\beta$ codifica la estabilidad (la cota es pequeña si $x(t_0)$ lo es) y la segunda la convergencia (la cota decae con el tiempo transcurrido $t-t_0$, no con el reloj $t$).
\end{ideaclave}

Volviendo al \cref{ej:vt-motivador}: para cada $t_0$ la solución cumple $\abs{x(t)}=\beta_{t_0}(\abs{x_0},t-t_0)$ con
\[
  \beta_{t_0}(r,s)=\frac{r(1+t_0)}{1+t_0+s}\in\KL .
\]
Pero $\sup_{t_0\ge0}\beta_{t_0}(r,s)=r$ para todo $s$: cualquier cota común tendría que valer al menos $r$ para todo $s$, y no decaería. No existe $\beta$ independiente de $t_0$ y el origen no es UAS (\cref{fig:vt-tiempo}, derecha).

\begin{definicion}[label=def:vt-exponencial]{Estabilidad exponencial}
El origen de \eqref{eq:vt-sistema} es \clave{exponencialmente estable} si existen $c,k,\lambda>0$ tales que
\[
  \norm{x(t)}\le k\,\norm{x(t_0)}\,e^{-\lambda(t-t_0)}\qquad\text{para todo }t\ge t_0\ge0\text{ y }\norm{x(t_0)}<c ,
\]
y \clave{globalmente exponencialmente estable} si esto vale para toda $x(t_0)$.
\end{definicion}

Es el caso particular $\beta(r,s)=kre^{-\lambda s}$, y por eso es automáticamente uniforme: en este curso, ``exponencialmente estable'' para un sistema variante quiere decir siempre con $k$ y $\lambda$ independientes de $t_0$.

%% ---------------------------------------------------------------------
\section{El método directo con \texorpdfstring{$V(t,x)$}{V(t,x)}}
\label{sec:vt-lyapunov}
%% ---------------------------------------------------------------------

Con un sistema variante en el tiempo es natural, y a veces necesario, que la función de Lyapunov también dependa de $t$. Hay que decidir entonces qué significa que $V(t,x)$ sea ``positiva'', y aparece una condición nueva, sin análogo en el \cref{cap:lyapunov}: que $V$ no crezca sin cota con el tiempo.

\subsection{Positividad, decrecimiento y radialidad}

\begin{definicion}[label=def:vt-decreciente]{Funciones $V(t,x)$ \citep[\S4.5]{Kha02}}
Sea $V:[0,\infty)\times\D\to\R$ continua con $V(t,0)=0$ para todo $t$. Se dice que $V$ es
\begin{itemize}
  \item \clave{definida positiva} si existe $W_1:\D\to\R$ continua y definida positiva con $V(t,x)\ge W_1(x)$ para todo $t\ge0$ y $x\in\D$;
  \item \clave{decreciente} si existe $W_2:\D\to\R$ continua y definida positiva con $V(t,x)\le W_2(x)$ para todo $t\ge0$ y $x\in\D$;
  \item \clave{radialmente no acotada} (con $\D=\R^n$) si $W_1$ se puede escoger radialmente no acotada.
\end{itemize}
Por el \cref{lem:vt-sandwich}, las dos primeras condiciones juntas equivalen a $\alpha_1(\norm x)\le V(t,x)\le\alpha_2(\norm x)$ con $\alpha_1,\alpha_2\in\K$ (en $\Kinf$ si es radialmente no acotada).
\end{definicion}

\notamargen{``Decreciente'' es la traducción usual del inglés \emph{decrescent}, y se presta a confusión: no dice que $V$ decrezca a lo largo de las soluciones (eso lo dice $\Vdot\le0$), sino que $V$ está acotada por arriba por una función de $x$ sola. \citet{KeBr23} usan el mismo término.}

Las dos cotas no tienen por qué depender del tiempo. ``Definida positiva'' prohíbe que $V(t,x)$ se acerque a cero con el tiempo para un $x\neq0$ fijo; ``decreciente'' prohíbe que se vaya a infinito. Ejemplos con $x\in\R^2$:

\begin{center}
\small
\begin{tabular}{@{}lccl@{}}
\toprule
$V(t,x)$ & d.\,p. & decreciente & razón \\
\midrule
$x_1^2(1+\sin^2t)+x_2^2(1+\cos^2t)$ & sí & sí & $\norm x^2\le V\le2\norm x^2$ \\
$x_1^2+(1+t)\,x_2^2$ & sí & no & en $x=(0,1)$ vale $1+t$ \\
$e^{-t}\norm x^2$ & no & sí & para $x$ fijo, $V\to0$ \\
$(2+\sin t)\,x^2$ & sí & sí & $x^2\le V\le3x^2$ \\
$x_1^2+x_1x_2^2(1+e^{-t})$ & no & sí & en $(-\epsilon^2,\epsilon)$ vale $-e^{-t}\epsilon^4<0$ \\
\bottomrule
\end{tabular}
\end{center}

Las dos primeras filas son el ejemplo~2.7 de \citet{KeBr23}. En la última, $V$ toma valores negativos en puntos arbitrariamente cercanos al origen: no es definida positiva en ningún $\Bola r$, aunque $x_1^2$ ``domine'' cuando $x_2=0$.

\subsection{La derivada a lo largo de las soluciones}

Si $V$ es de clase $\Cont1$, a lo largo de una solución de \eqref{eq:vt-sistema},
\begin{equation}
  \Vdot(t,x)=\frac{\partial V}{\partial t}(t,x)+\frac{\partial V}{\partial x}(t,x)\,f(t,x) .
  \label{eq:vt-vdot}
\end{equation}
El primer término es nuevo: mide cuánto cambia $V$ por el paso del reloj, con $x$ congelado.

\begin{trampa}[label=tr:vt-parcial]{``$\Vdot=\nabla V\cdot f$, como siempre''}
\plausible{Es la fórmula de los siete capítulos anteriores.}
Para $\dot x=-x/10$ y $V=(2+\sin t)\,x^2$, olvidando $\partial V/\partial t$ se obtiene $\nabla V\cdot f=-0.2(2+\sin t)\,x^2<0$ para todo $x\neq0$. La derivada correcta es
\[
  \Vdot=\cos t\;x^2-0.2(2+\sin t)\,x^2=\bigl(\cos t-0.2(2+\sin t)\bigr)x^2 ,
\]
que en $t=0$ vale $0.6\,x^2>0$: $V$ \emph{sube} a lo largo de la solución durante parte de cada periodo (\cref{fig:vt-parcial}). Aquí el sistema es exponencialmente estable de todos modos (basta $V=x^2$), pero la conclusión ``$\Vdot<0$'' era falsa.
\correcto{Use siempre \eqref{eq:vt-vdot}. Si $V$ no depende de $t$, el término desaparece; si depende, casi siempre es el que decide.}
\end{trampa}

\begin{figure}[ht]
\centering
\includegraphics[width=\linewidth]{cap09/f_derivada_parcial}
\caption{Izquierda: $V(t,x(t))=(2+\sin t)\,x(t)^2$ a lo largo de una solución de $\dot x=-x/10$: baja a la larga, pero sube a ratos. Derecha: el coeficiente de $x^2$ en $\Vdot$ con el término $\partial V/\partial t$ (rojo) y sin él (gris punteado).}
\label{fig:vt-parcial}
\end{figure}

\codigorepo{cap09_variantes/C9_Ej3_DerivadaParcial.m}{derivada}{$\Vdot$ con y sin $\partial V/\partial t$, comparado con la derivada numérica de $V(t,x(t))$}

\subsection{Los teoremas}

\begin{teorema}[label=thm:vt-lyapunov]{Lyapunov para sistemas variantes \citep[Teo.~4.8 y 4.9]{Kha02}}
Sea $V:[0,\infty)\times\D\to\R$ de clase $\Cont1$ con
\begin{equation}
  W_1(x)\le V(t,x)\le W_2(x)\qquad\text{para todo }t\ge0,\ x\in\D ,
  \label{eq:vt-w1w2}
\end{equation}
con $W_1,W_2$ continuas y definidas positivas en $\D$.
\begin{enumerate}
  \item Si $\Vdot(t,x)\le0$ para todo $t\ge0$ y $x\in\D$, el origen es \clave{uniformemente estable}.
  \item Si $\Vdot(t,x)\le-W_3(x)$ para todo $t\ge0$ y $x\in\D$, con $W_3$ continua y definida positiva, el origen es \clave{uniformemente asintóticamente estable}.
  \item Si además $\D=\R^n$ y $W_1$ es radialmente no acotada, el origen es \clave{globalmente uniformemente asintóticamente estable}.
\end{enumerate}
\end{teorema}

\begin{demostracion}
Se prueba (2) con funciones de comparación; (1) es el primer paso y (3) es el mismo argumento con $\Kinf$. Por el \cref{lem:vt-sandwich} hay $\alpha_1,\alpha_2,\alpha_3\in\K$ en $[0,r)$, con $\Bola r\subset\D$, tales que $\alpha_1(\norm x)\le V(t,x)\le\alpha_2(\norm x)$ y $W_3(x)\ge\alpha_3(\norm x)$; sin pérdida de generalidad, $\alpha_3$ es localmente Lipschitz.

\emph{Paso 1 (las soluciones no salen).} Sea $\rho<\alpha_1(r)$ y $c>0$ con $\alpha_2(c)<\rho$. Si $\norm{x(t_0)}<c$, entonces $V(t_0,x(t_0))<\rho$, y como $V$ no crece, $\alpha_1(\norm{x(t)})\le V(t,x(t))<\rho$: la solución se queda en $\Bola{\alpha_1^{-1}(\rho)}\subset\Bola r$ para todo $t\ge t_0$. Con $\Vdot\le0$ esto ya da $\norm{x(t)}\le\alpha_1^{-1}(\alpha_2(\norm{x(t_0)}))$, que es \eqref{eq:vt-us} con una función de $\K$ que no depende de $t_0$: US.

\emph{Paso 2 (desigualdad diferencial escalar).} Con $\norm x\ge\alpha_2^{-1}(V)$,
\[
  \Vdot\le-\alpha_3(\norm x)\le-\alpha_3\bigl(\alpha_2^{-1}(V)\bigr)=:-\alpha(V) ,\qquad\alpha=\alpha_3\circ\alpha_2^{-1}\in\K .
\]

\emph{Paso 3 (comparación).} Por el \cref{lem:vt-flujo}, $\dot y=-\alpha(y)$ tiene solución $\sigma(y_0,t-t_0)$ con $\sigma\in\KL$, y por el lema de comparación (\cref{lem:pm-comparacion}), $V(t,x(t))\le\sigma\bigl(V(t_0,x(t_0)),t-t_0\bigr)$.

\emph{Paso 4 (de vuelta a $\norm x$).} Con $\alpha_1(\norm x)\le V$ y $V(t_0,x_0)\le\alpha_2(\norm{x_0})$,
\[
  \norm{x(t)}\le\alpha_1^{-1}\Bigl(\sigma\bigl(\alpha_2(\norm{x(t_0)}),\,t-t_0\bigr)\Bigr)=:\beta\bigl(\norm{x(t_0)},t-t_0\bigr) ,
\]
y $\beta\in\KL$ por el \cref{lem:vt-propiedades}. Ninguna de las funciones depende de $t_0$: por el \cref{lem:vt-kl}, el origen es UAS. En el caso (3), $\alpha_1,\alpha_2\in\Kinf$, el paso~1 vale para todo $c$ y la cota vale en todo $\R^n$.
\end{demostracion}

\notamargen{Dónde entra cada hipótesis: $W_1$ (definida positiva) en el paso~4, para pasar de $V$ a $\norm x$; $W_2$ (decreciente) en los pasos~1 y~2, para que $c$ y $\alpha$ no dependan de $t_0$; $W_3$ en el paso~2. Sin $W_2$ el argumento da estabilidad y convergencia, pero no uniformes.}

\begin{teorema}[label=thm:vt-exp]{Estabilidad exponencial \citep[Teo.~4.10]{Kha02}}
Si existen $k_1,k_2,k_3,a>0$ tales que
\[
  k_1\norm x^a\le V(t,x)\le k_2\norm x^a ,\qquad\Vdot(t,x)\le-k_3\norm x^a
\]
para todo $t\ge0$ y $x\in\D$, el origen es exponencialmente estable, con
\begin{equation}
  \norm{x(t)}\le\Bigl(\frac{k_2}{k_1}\Bigr)^{1/a}\norm{x(t_0)}\,e^{-\frac{k_3}{a\,k_2}(t-t_0)} .
  \label{eq:vt-cotaexp}
\end{equation}
Si $\D=\R^n$, es globalmente exponencialmente estable.
\end{teorema}

\begin{demostracion}
$\Vdot\le-k_3\norm x^a\le-\frac{k_3}{k_2}V$, así que $V(t)\le V(t_0)e^{-\frac{k_3}{k_2}(t-t_0)}$ (\cref{lem:pm-comparacion}). Entonces $k_1\norm{x(t)}^a\le k_2\norm{x(t_0)}^ae^{-\frac{k_3}{k_2}(t-t_0)}$, y se saca raíz $a$-ésima.
\end{demostracion}

\begin{center}
\small
\begin{tabular}{@{}p{4.9cm}p{3.1cm}p{3.3cm}@{}}
\toprule
Hipótesis sobre $V(t,x)$ & $\Vdot$ & Conclusión \\
\midrule
definida positiva & $\le0$ & estable \\
definida positiva y decreciente & $\le0$ & US \\
definida positiva y decreciente & $\le-W_3(x)$ & UAS \\
\quad$+$ radialmente no acotada, $\D=\R^n$ & $\le-W_3(x)$ & GUAS \\
$k_1\norm x^a\le V\le k_2\norm x^a$ & $\le-k_3\norm x^a$ & exponencial (uniforme) \\
\bottomrule
\end{tabular}
\end{center}

La primera fila, estabilidad no uniforme con $V$ definida positiva no decreciente, es la que demostró Lyapunov para el caso variante~\citep{Lya92}; en la práctica casi no se usa, porque sin uniformidad no hay robustez. Si $V$ no depende de $t$, las dos cotas \eqref{eq:vt-w1w2} se cumplen con $W_1=W_2=V$, y el teorema se reduce al del \cref{cap:lyapunov}: en ese caso, la uniformidad es gratis.

\subsection{La trampa de $\Vdot<0$}

\begin{ejemplo}[label=ej:vt-e1]{$\Vdot<0$ y aun así no hay convergencia a cero}
Para $\dot x=-\dfrac{x}{(1+t)^2}$ con $V=x^2$, decida si el origen es asintóticamente estable.
\solucion $V$ es definida positiva, decreciente y radialmente no acotada, y
\[
  \Vdot=-\frac{2x^2}{(1+t)^2}<0\qquad\text{para todo }x\neq0 .
\]
Pero integrando,
\[
  x(t)=x_0\exp\Bigl(-\int_{t_0}^t\frac{\dif\tau}{(1+\tau)^2}\Bigr)=x_0\,e^{\frac1{1+t}-\frac1{1+t_0}}\xrightarrow[t\to\infty]{}x_0\,e^{-\frac1{1+t_0}}\neq0 .
\]
Con $x_0=1$, el límite es $e^{-1}=0.368$ si $t_0=0$ y $e^{-1/11}=0.913$ si $t_0=10$ (\cref{fig:vt-e1}). El origen es uniformemente estable ($\abs{x(t)}\le\abs{x_0}$), pero no es asintóticamente estable.
\end{ejemplo}

\begin{trampa}[label=tr:vt-negativa]{``$V>0$ y $\Vdot<0$ fuera del origen, luego AS''}
\plausible{Para $\dot x=f(x)$, $V$ definida positiva y $\Vdot$ definida negativa dan estabilidad asintótica (\cref{thm:ly-directo}).}
En el \cref{ej:vt-e1}, $\Vdot<0$ para todo $x\neq0$, pero para $x$ fijo $-\Vdot=2x^2/(1+t)^2\to0$: el decrecimiento se apaga con el reloj, y ninguna $W_3(x)$ definida positiva queda por debajo de $-\Vdot$ para todo $t$. La energía que se disipa en total es finita y no alcanza para llegar a cero.
\correcto{Lo que se necesita es $\Vdot\le-W_3(x)$, con $W_3$ \emph{independiente de $t$}. ``$\Vdot<0$ para cada $t$'' no es lo mismo que ``$\Vdot$ definida negativa'' en el sentido de la \cref{def:vt-decreciente}.}
\end{trampa}

\begin{figure}[ht]
\centering
\includegraphics[width=0.75\linewidth]{cap09/f_trampa_e1}
\caption{$\dot x=-x/(1+t)^2$ con $x_0=1$ y varios $t_0$: cada solución se estanca en $e^{-1/(1+t_0)}$ (punteadas), aunque $\Vdot<0$ en todo instante.}
\label{fig:vt-e1}
\end{figure}

\codigorepo{cap09_variantes/C9_Ej4_TrampaE1.m}{trampae1}{Límites $x_0e^{-1/(1+t_0)}$ contra la simulación}

%% ---------------------------------------------------------------------
\section{Ejemplos y sistemas lineales variantes}
\label{sec:vt-ejemplos}
%% ---------------------------------------------------------------------

\subsection{La ecuación de Mathieu amortiguada}

\begin{ejemplo}[label=ej:vt-mathieu]{Un resorte cuya rigidez oscila}
Considere $\ddot y+\dot y+(2+\sin t)\,y=0$, es decir,
\begin{equation}
  \dot x_1=x_2,\qquad\dot x_2=-x_2-(2+\sin t)\,x_1 ,
  \label{eq:vt-mathieu}
\end{equation}
y la candidata $V(t,x)=x_1^2+\dfrac{x_2^2}{2+\sin t}$. ¿Qué concluye el \cref{thm:vt-lyapunov}?
\solucion Es la energía con la rigidez del instante: si la rigidez fuera constante $k$, $kx_1^2+x_2^2$ sería la energía, y dividir por $k$ da $V$. Como $1\le2+\sin t\le3$,
\[
  x_1^2+\tfrac13x_2^2\le V(t,x)\le x_1^2+x_2^2 ,
\]
así que $V$ es definida positiva, decreciente y radialmente no acotada. Con \eqref{eq:vt-vdot},
\[
  \frac{\partial V}{\partial t}=-\frac{\cos t\;x_2^2}{(2+\sin t)^2},\qquad
  \frac{\partial V}{\partial x}f=2x_1x_2+\frac{2x_2}{2+\sin t}\bigl(-x_2-(2+\sin t)x_1\bigr)=-\frac{2x_2^2}{2+\sin t} ,
\]
y sumando,
\begin{equation}
  \Vdot=-\frac{4+2\sin t+\cos t}{(2+\sin t)^2}\,x_2^2\le-c\,x_2^2 .
  \label{eq:vt-mathieuvdot}
\end{equation}
El numerador es al menos $4-\sqrt5=1.76$ (porque $2\sin t+\cos t\ge-\sqrt5$) y el denominador a lo sumo $9$, de modo que $c=(4-\sqrt5)/9=0.196$ sirve; el mínimo exacto del cociente es $0.637$. $\Vdot$ es semidefinida negativa: se anula en todo el eje $x_2=0$. El teorema da \clave{estabilidad uniforme global}, y nada más.
\end{ejemplo}

\begin{figure}[ht]
\centering
\includegraphics[width=\linewidth]{cap09/f_mathieu}
\caption{Mathieu amortiguada desde $(1.5,0)$. Izquierda: la solución y el conjunto de nivel $V(t,x)=2.25$ en tres instantes; la elipse ``respira'' con $t$, siempre entre la circunferencia $x_1^2+x_2^2=2.25$ ($t=4.71$) y la elipse $x_1^2+\tfrac13x_2^2=2.25$ ($t=1.57$), que son los conjuntos de nivel de las cotas $W_2$ y $W_1$ de \eqref{eq:vt-w1w2}. Derecha: $V(t,x(t))$ baja con mesetas (donde $x_2\approx0$) y, de hecho, decae como $e^{-t}$.}
\label{fig:vt-mathieu}
\end{figure}

La figura sugiere que el origen es mucho mejor que uniformemente estable. Como \eqref{eq:vt-mathieu} es lineal y periódico (periodo $2\pi$), la teoría de Floquet lo decide. Para $\dot x=A(t)x$ con $A(t+\tau)=A(t)$, la matriz de monodromía $M=\Phi(\tau,0)$ es el mapa estroboscópico de la \cref{sec:bf-poincare} (aquí lineal): $x(k\tau)=M^kx(0)$. Sus valores propios, los \clave{multiplicadores de Floquet}, deciden: si todos tienen módulo menor que $1$, el origen es uniforme y exponencialmente estable; si alguno tiene módulo mayor que $1$, es inestable \citep[\S7.3, Prop.~7.14]{Sas99}. Para \eqref{eq:vt-mathieu}, $\tau=2\pi$ y $M$, calculada numéricamente, tiene los multiplicadores
\[
  \mu=-0.0099\pm0.0421\,j,\qquad\abs\mu=e^{-\pi}=0.0432<1 .
\]
El módulo no es casualidad: por la fórmula de Liouville, $\det M=\exp\int_0^{2\pi}\operatorname{tr}A(\tau)\dif\tau=e^{-2\pi}$, y como los multiplicadores son complejos conjugados, $\abs\mu^2=e^{-2\pi}$. Con los dos multiplicadores dentro del círculo unitario el origen es \clave{uniforme y globalmente exponencialmente estable}, con tasa $\ln\abs\mu/2\pi=-\tfrac12$: $\norm{x(t)}e^{t/2}$ se mantiene acotado (por $2.05$ en $[10,40]$ desde $(1,0)$). En la \cref{sec:vt-semidefinida} se cierra la brecha sin calcular $M$, con LaSalle--Yoshizawa y Barbalat.

\codigorepo{cap09_variantes/C9_Ej5_Mathieu.m}{mathieu}{$\Vdot$ simbólica, la cota $c$ y los multiplicadores de Floquet}

\video{V9_2_mathieu}{Los niveles de $V(t,x)$ respiran y la solución entra}

\subsection{Una ganancia variable}

\begin{ejemplo}[label=ej:vt-ganancia]{GES con una condición sobre $\dot g$ \citep[\S4.5]{Kha02}}
Considere
\[
  \dot x_1=-x_1-g(t)\,x_2,\qquad\dot x_2=x_1-x_2 ,
\]
con $g$ de clase $\Cont1$, $0\le g(t)\le k$ y $\dot g(t)\le g(t)$. Pruebe que el origen es globalmente exponencialmente estable.
\solucion Tome $V=x_1^2+(1+g(t))\,x_2^2$, de modo que $\norm x^2\le V\le(1+k)\norm x^2$. Entonces
\[
  \Vdot=\dot g\,x_2^2+2x_1(-x_1-gx_2)+2(1+g)x_2(x_1-x_2)=-2x_1^2+2x_1x_2-\bigl(2+2g-\dot g\bigr)x_2^2 .
\]
Los términos cruzados $\mp2gx_1x_2$ se cancelan: para eso está el factor $(1+g)$. Con $\dot g\le g$, $2+2g-\dot g\ge2+g\ge2$, y
\[
  \Vdot\le-2x_1^2+2x_1x_2-2x_2^2=-x\T\begin{bmatrix}2&-1\\-1&2\end{bmatrix}x\le-\norm x^2 ,
\]
porque los valores propios de la matriz son $1$ y $3$. El \cref{thm:vt-exp} con $k_1=1$, $k_2=1+k$, $k_3=1$, $a=2$ da
\[
  \norm{x(t)}\le\sqrt{1+k}\;\norm{x(t_0)}\;e^{-\frac{t-t_0}{2(1+k)}} .
\]
\end{ejemplo}

Con $g=1+0.5\sin t$ ($k=1.5$), $g-\dot g\ge1-0.5\sqrt2=0.29>0$ y la cota se cumple con holgura: en la simulación de la \cref{fig:vt-ganancia}, el cociente entre $\norm{x(t)}$ y la cota vale $1/\sqrt{1+k}=0.63$ en $t=t_0$ y después solo baja. Con $g=(1+\sin10t)/2$, en cambio, $2+2g-\dot g$ vale $-2$ en $t=0$: el término $\partial V/\partial t$ inyecta más de lo que el sistema disipa y $\Vdot$ deja de ser negativa. \alerta{Eso no prueba que el sistema sea inestable}: solo que esta $V$ ya no es un certificado.

\begin{figure}[ht]
\centering
\includegraphics[width=\linewidth]{cap09/f_ganancia}
\caption{Izquierda: $\norm{x(t)}$ para varias condiciones iniciales en el círculo unitario con $g=1+0.5\sin t$, y la cota $\sqrt{1+k}\,e^{-(t-t_0)/(2(1+k))}$. Derecha: el coeficiente $2+2g-\dot g$ de $x_2^2$; con la $g$ lenta se queda sobre $2$, con la rápida cruza a valores negativos.}
\label{fig:vt-ganancia}
\end{figure}

\codigorepo{cap09_variantes/C9_Ej6_GananciaVariable.m}{ganancia}{La cota exponencial y el coeficiente $2+2g-\dot g$}

\subsection{Sistemas lineales variantes: los valores propios no bastan}

Para $\dot x=A(t)x$ la solución es $x(t)=\Phi(t,t_0)x(t_0)$, con $\Phi$ la matriz de transición. Las nociones de la \cref{sec:vt-comparacion} se simplifican mucho.

\begin{teorema}[label=thm:vt-lineal]{Estabilidad de $\dot x=A(t)x$ \citep[Teo.~4.11 y 4.12]{Kha02}}
\begin{enumerate}
  \item El origen es UAS si y solo si es globalmente exponencialmente estable, si y solo si existen $k,\lambda>0$ con $\norm{\Phi(t,t_0)}\le ke^{-\lambda(t-t_0)}$ para todo $t\ge t_0\ge0$.
  \item Si existen $P(t)$ de clase $\Cont1$, simétrica, y constantes $c_1,c_2,c_3>0$ con $c_1I\le P(t)\le c_2I$ y
  \begin{equation}
    \dot P(t)+A\T(t)P(t)+P(t)A(t)=-Q(t),\qquad Q(t)\ge c_3I ,
    \label{eq:vt-lyapdif}
  \end{equation}
  el origen es globalmente exponencialmente estable. Recíprocamente, si $A(t)$ es continua y acotada y el origen es UAS, para cada $Q(t)$ continua, acotada y con $Q\ge c_3I$, la matriz $P(t)=\int_t^\infty\Phi\T(\tau,t)Q(\tau)\Phi(\tau,t)\dif\tau$ cumple \eqref{eq:vt-lyapdif} y las cotas.
\end{enumerate}
\end{teorema}

La parte de suficiencia es el \cref{thm:vt-exp} con $V=x\T P(t)x$: $\Vdot=x\T(\dot P+A\T P+PA)x=-x\T Qx\le-c_3\norm x^2$. La ecuación \eqref{eq:vt-lyapdif} es la ecuación de Lyapunov \emph{diferencial}: para $A$ constante y $P$ constante es la del \cref{thm:ly-ecuacion}. Lo que \emph{no} sirve es congelar el tiempo.

\begin{ejemplo}[label=ej:vt-markus]{Markus--Yamabe \citep{MaYa60}, \citep[\S4.6]{Kha02}}
Considere $\dot x=A(t)x$ con
\[
  A(t)=\begin{bmatrix}-1+1.5\cos^2t & 1-1.5\sin t\cos t\\[2pt] -1-1.5\sin t\cos t & -1+1.5\sin^2t\end{bmatrix} .
\]
Calcule los valores propios de $A(t)$ y verifique que
\[
  \Phi(t,0)=\begin{bmatrix}e^{t/2}\cos t & e^{-t}\sin t\\ -e^{t/2}\sin t & e^{-t}\cos t\end{bmatrix} .
\]
\solucion Con $c=\cos t$ y $s=\sin t$, $\operatorname{tr}A=-2+1.5(c^2+s^2)=-\tfrac12$ y $\det A=(-1+1.5c^2)(-1+1.5s^2)-(1-1.5sc)(-1-1.5sc)=\tfrac12$. El polinomio característico $\lambda^2+\tfrac12\lambda+\tfrac12$ no depende de $t$:
\[
  \lambda=-\tfrac14\pm j\,\tfrac{\sqrt7}4\qquad\text{para todo }t .
\]
Para la primera columna $\varphi=e^{t/2}(c,-s)$: $\dot\varphi=e^{t/2}\bigl(\tfrac12c-s,\ -\tfrac12s-c\bigr)$, y $A\varphi=e^{t/2}\bigl(-c+1.5c^3-s+1.5s^2c,\ -c-1.5sc^2+s-1.5s^3\bigr)=e^{t/2}\bigl(\tfrac12c-s,\ -c-\tfrac12s\bigr)$. La segunda columna se verifica igual, y $\det\Phi=e^{-t/2}=\exp\int_0^t\operatorname{tr}A$, como pide Liouville. Desde $x(0)=(1,0)$, $x(t)=e^{t/2}(\cos t,-\sin t)$ y $\norm{x(t)}=e^{t/2}$: \clave{el origen es inestable}.
\end{ejemplo}

\begin{trampa}[label=tr:vt-congelar]{``Si $A(t)$ es Hurwitz para cada $t$, el sistema es estable''}
\plausible{Para cada instante, el sistema ``congelado'' $\dot x=A(t^*)x$ es asintóticamente estable, con los mismos valores propios.}
En el \cref{ej:vt-markus} los valores propios congelados son $-\tfrac14\pm j\tfrac{\sqrt7}4$ para todo $t$, y las soluciones crecen como $e^{t/2}$. La factorización $\Phi(t,0)=R(-t)\,\diag(e^{t/2},e^{-t})$, con $R(\theta)$ la rotación de ángulo $\theta$, explica lo que pasa: en un marco que gira con el campo, el sistema es una silla con valor propio $+\tfrac12$, y la rotación lo esconde de los valores propios (\cref{fig:vt-markus}).
\correcto{Los valores propios de $A(t)$ no deciden la estabilidad de $\dot x=A(t)x$. Sirven la matriz de transición (\cref{thm:vt-lineal}), la ecuación de Lyapunov diferencial, Floquet en el caso periódico, o resultados de variación lenta: si $A(t)$ es Hurwitz uniformemente y $\norm{\dot A(t)}$ es suficientemente pequeña, el origen sí es exponencialmente estable \citep[\S9.6]{Kha02}.}
\end{trampa}

\begin{figure}[ht]
\centering
\includegraphics[width=\linewidth]{cap09/f_markus}
\caption{Markus--Yamabe desde $(1,0)$. Izquierda: la solución real (coral) se va en espiral; las de los campos congelados $\dot x=A(t^*)x$ (punteadas) convergen al origen. Derecha: $\norm{x(t)}$ coincide con $e^{t/2}$, mientras los valores propios ``dicen'' $e^{-t/4}$.}
\label{fig:vt-markus}
\end{figure}

\codigorepo{cap09_variantes/C9_Ej7_MarkusYamabe.m}{markus}{Valores propios congelados contra la matriz de transición}

\video{V9_3_markus_yamabe}{Los campos congelados atraen; la solución real se escapa}

\subsection{La linealización necesita uniformidad}

El método indirecto (\cref{thm:ly-indirecto}) tiene una versión variante \citep[Teo.~4.13]{Kha02}: si $f(t,x)$ es de clase $\Cont1$ y su jacobiano $\partial f/\partial x$ es \emph{acotado y Lipschitz en $\Bola r$, uniformemente en $t$}, y el origen de la linealización $\dot x=A(t)x$, con $A(t)=\frac{\partial f}{\partial x}(t,0)$, es exponencialmente estable, entonces el origen del sistema no lineal es exponencialmente estable. La hipótesis de uniformidad no es decorativa.

\begin{ejemplo}[label=ej:vt-linealizacion]{Una linealización Hurwitz que engaña}
Para $\dot x_1=-x_1+t\,x_2^2$, $\dot x_2=x_1-x_2$, la linealización en el origen es $\dot x=\bigl[\begin{smallmatrix}-1&0\\1&-1\end{smallmatrix}\bigr]x$, constante y Hurwitz. Pero el jacobiano $\bigl[\begin{smallmatrix}-1&2tx_2\\1&-1\end{smallmatrix}\bigr]$ no es Lipschitz uniformemente en $t$: el término cuadrático tiene un coeficiente que crece con el reloj. Desde la misma condición inicial $(0.1,0.1)$, la solución converge si $t_0=0$ y explota si $t_0=10$ (cruza $\norm x=10^3$ en $t=16.33$) o $t_0=100$ (\cref{fig:vt-linealizacion}). El radio de la región que converge depende de $t_0$ y se encoge: no hay un $c$ uniforme, y el origen no es UAS.
\end{ejemplo}

\begin{figure}[ht]
\centering
\includegraphics[width=0.8\linewidth]{cap09/f_linealizacion}
\caption{$\norm{x}$ contra $t-t_0$ desde $(0.1,0.1)$ para $\dot x_1=-x_1+tx_2^2$, $\dot x_2=x_1-x_2$. La linealización (gris, a trazos) es la misma para todo $t_0$; el sistema no lineal converge con $t_0=0$ y explota con $t_0=10$ y $t_0=100$.}
\label{fig:vt-linealizacion}
\end{figure}

\codigorepo{cap09_variantes/C9_Ej8_Linealizacion.m}{linealizacion}{La misma condición inicial con tres horas de arranque}

%% ---------------------------------------------------------------------
\section{Cuando \texorpdfstring{$\Vdot$}{dV/dt} es solo semidefinida}
\label{sec:vt-semidefinida}
%% ---------------------------------------------------------------------

En la Mathieu amortiguada (\cref{ej:vt-mathieu}), $\Vdot\le-c\,x_2^2$: semidefinida negativa. Para un sistema invariante en el tiempo se usaría LaSalle (\cref{thm:iv-lasalle}), pero su prueba descansa en que los conjuntos límite son invariantes, y eso es falso cuando $f$ depende de $t$: una solución que se acerca a un conjunto no tiene por qué ``ver'' el mismo campo vectorial cuando llega. Las herramientas disponibles son estas.

\begin{center}
\small
\begin{tabular}{@{}p{3.6cm}p{3.7cm}p{4.6cm}@{}}
\toprule
Sistema & Herramienta & Qué da \\
\midrule
$\dot x=f(x)$ & LaSalle (\cref{thm:iv-lasalle}) & convergencia al mayor invariante en $\{\Vdot=0\}$ \\
$f(t+\tau,x)=f(t,x)$ & LaSalle periódico & lo mismo, con el mapa estroboscópico \\
$\dot x=f(t,x)$ general & LaSalle--Yoshizawa (\cref{thm:vt-ly}) & $W(x(t))\to0$ si $\Vdot\le-W(x)\le0$ \\
$\dot x=f(t,x)$ general & Barbalat (\cref{lem:vt-barbalat}) & completar la convergencia en las variables que faltan \\
$\dot x=f(t,x)$ general & Matrosov & UAS con varias funciones auxiliares \\
\bottomrule
\end{tabular}
\end{center}

LaSalle--Yoshizawa \citep{Yos66} es la generalización directa: con $V$ definida positiva y decreciente y $\Vdot\le-W(x)\le0$, las soluciones están acotadas y $W(x(t))\to0$. Da convergencia de la parte de $x$ que $W$ ``ve'', no de todo el estado. Su prueba descansa en un resultado sobre señales, el lema de Barbalat, y los dos se prueban a continuación.

\subsection{Señales: el lema de Barbalat}

Aquí la discusión cambia de escenario: \alerta{el lema de Barbalat es un resultado sobre funciones de una variable real (señales), no sobre ecuaciones diferenciales}. Se aplica a la señal $t\mapsto V(t,x(t))$, o a $t\mapsto W(x(t))$, sobre una solución fija. La hipótesis vive en el espacio de estados; la conclusión, en el eje del tiempo.

En la situación de esta sección, $\Vdot(t,x)\le-W(x)\le0$ con $W$ solo semidefinida positiva, hay algo que sí es gratis: si $V$ está acotada por debajo y no crece a lo largo de la solución, converge a un límite $V_\infty$, e integrando
\begin{equation}
  \int_{t_0}^\infty W\bigl(x(t)\bigr)\dif t\le V\bigl(t_0,x(t_0)\bigr)-V_\infty<\infty .
  \label{eq:vt-integral}
\end{equation}

\begin{trampa}[label=tr:vt-integral]{``La integral es finita, luego $W(x(t))\to0$''}
\plausible{Si el área bajo una curva no negativa es finita, la curva tiene que irse a cero.}
Una función no negativa puede tener integral finita sin tender a cero: basta con pulsos cada vez más angostos. Con pulsos triangulares centrados en $t=k$, de altura $k$ y base $1/k^3$, el área total es $\sum_k\frac1{2k^2}=\pi^2/12$, pero $\psi(k)=k\to\infty$ (\cref{fig:vt-pulsos}).
\correcto{Falta una hipótesis de regularidad: la continuidad uniforme. Los pulsos tienen pendientes $\pm2k^4$, no acotadas.}
\end{trampa}

\begin{figure}[ht]
\centering
\includegraphics[width=\linewidth]{cap09/f_pulsos}
\caption{Tres trampas sobre señales. Izquierda: pulsos de área total $\pi^2/12$ con alturas que crecen sin cota. Centro: $h=\sin(\ln t)$ tiene $\dot h\to0$ pero no converge. Derecha: $h=e^{-t}\sin(e^{2t})\to0$ pero $\dot h$ no está acotada.}
\label{fig:vt-pulsos}
\end{figure}

Las otras dos trampas de la figura: $\dot h\to0$ no implica que $h$ converja ($h=\sin(\ln t)$ oscila entre $\pm1$ cada vez más despacio), y $h$ convergente no implica $\dot h\to0$ ($h=e^{-t}\sin(e^{2t})$: la amplitud se apaga, pero la frecuencia crece más rápido).

\begin{definicion}[label=def:vt-unifcont]{Continuidad uniforme}
$\psi:[0,\infty)\to\R$ es \clave{uniformemente continua} si para todo $\varepsilon>0$ existe $\delta(\varepsilon)>0$, \emph{el mismo para todo $t$}, tal que $\abs{t_1-t_2}<\delta\implica\abs{\psi(t_1)-\psi(t_2)}<\varepsilon$.
\end{definicion}

En la práctica casi siempre se verifica así: \alerta{si $\abs{\dot\psi(t)}\le L$ para todo $t$, entonces $\psi$ es uniformemente continua} (valor medio, $\delta=\varepsilon/L$). $\sin t$, $e^{-t}$ y $\arctan t$ lo son; $\sin(t^2)$ no (entre $t_n=\sqrt{2\pi n}$ y $s_n=\sqrt{2\pi n+\pi/2}$ la distancia tiende a cero y los valores difieren en $1$); los pulsos tampoco. Es la misma idea de ``uniforme'' de la \cref{sec:vt-uniformidad}: una constante que no depende del instante (\cref{def:pm-unifcont}).

\begin{lema}[label=lem:vt-barbalat]{Barbalat \citep{Bar59}; \citep[Lema~8.2]{Kha02}}
\textbf{Forma integral.} Si $\psi$ es uniformemente continua y $\lim_{t\to\infty}\int_0^t\psi(\tau)\dif\tau$ existe y es finito, entonces $\psi(t)\to0$.

\textbf{Forma de derivada.} Si $h$ tiene límite finito cuando $t\to\infty$ y $\dot h$ es uniformemente continua, entonces $\dot h(t)\to0$. En particular, basta con que $\ddot h$ exista y esté acotada.
\end{lema}

\begin{demostracion}
Si $\psi\not\to0$, hay $\varepsilon>0$ y $t_k\to\infty$ con $\abs{\psi(t_k)}\ge\varepsilon$. Con $\delta=\delta(\varepsilon/2)$, en $[t_k,t_k+\delta]$ se tiene $\abs{\psi}>\varepsilon/2$ y $\psi$ no cambia de signo, así que $\bigl|\int_{t_k}^{t_k+\delta}\psi\bigr|\ge\varepsilon\delta/2$ para todo $k$. Pero si $\int_0^t\psi$ converge, el criterio de Cauchy exige que esas integrales tiendan a cero. La clave es que $\delta$ es \emph{el mismo} para todo $k$: sin continuidad uniforme, los intervalos podrían encogerse, como en los pulsos. La forma de derivada es la integral con $\psi=\dot h$.
\end{demostracion}

\codigorepo{cap09_variantes/C9_Ej12_Senales.m}{pulsos}{Las tres trampas sobre señales}

\subsection{El teorema de LaSalle--Yoshizawa}

\begin{teorema}[label=thm:vt-ly]{LaSalle--Yoshizawa \citep{Yos66}; \citep[Teo.~8.4]{Kha02}}
Sea $\dot x=f(t,x)$ con $f$ continua a trozos en $t$, localmente Lipschitz en $x$ uniformemente en $t$, y $f(t,0)=0$. Si $V\in\Cont1$ cumple $\alpha_1(\norm x)\le V(t,x)\le\alpha_2(\norm x)$ y $\Vdot\le-W(x)\le0$ en $[0,\infty)\times\Bola r$, con $W$ continua, entonces para $\norm{x(t_0)}<\alpha_2^{-1}(\alpha_1(r))$ las soluciones están acotadas y $W(x(t))\to0$. Si las hipótesis valen en todo $\R^n$ y $\alpha_1\in\Kinf$, la conclusión vale para toda $x(t_0)$.
\end{teorema}

\begin{demostracion}
(i) $V$ no crece, luego $x(t)$ se queda en un compacto $\Bola\rho\subset\Bola r$, con $\rho=\alpha_1^{-1}(\alpha_2(\norm{x(t_0)}))$; en el caso global, $\rho$ es finito para toda $x(t_0)$. (ii) Vale \eqref{eq:vt-integral}. (iii) En $\Bola\rho$, $\norm{f}\le L\rho$: $x(t)$ es Lipschitz en $t$; $W$ es uniformemente continua en el compacto; la composición $t\mapsto W(x(t))$ es uniformemente continua. (iv) Barbalat, forma integral.
\end{demostracion}


\begin{ejemplo}[label=ej:vt-barbalat]{Cerrar la Mathieu amortiguada con Barbalat}
Con los datos del \cref{ej:vt-mathieu}, pruebe que $x(t)\to0$ para toda condición inicial.
\solucion
\emph{Paso 1 (LaSalle--Yoshizawa).} $V$ es definida positiva, decreciente y radialmente no acotada y $\Vdot\le-c\,x_2^2=:-W(x)$: las soluciones están acotadas y $x_2(t)\to0$.

\emph{Paso 2 ($\dot x_2$ es uniformemente continua).} Derivando la segunda ecuación,
\[
  \ddot x_2=-\dot x_2-\cos t\;x_1-(2+\sin t)\,\dot x_1 ,
\]
que está acotada porque $x$, $\dot x$, $\sin t$ y $\cos t$ lo están (desde $(1,0)$, $\abs{\ddot x_2}\le2.77$).

\emph{Paso 3 (la integral converge).} $\int_{t_0}^t\dot x_2(\tau)\dif\tau=x_2(t)-x_2(t_0)\to-x_2(t_0)$.

\emph{Paso 4 (Barbalat).} $\dot x_2(t)\to0$.

\emph{Paso 5 (la variable que faltaba).} De la ecuación, $(2+\sin t)\,x_1=-\dot x_2-x_2\to0$, y como $2+\sin t\ge1$, $x_1(t)\to0$.

Con el paso~1, el origen es uniformemente estable y globalmente atractivo. Como el sistema es periódico, la atracción es uniforme (\cref{sec:vt-uniformidad}), y como es lineal, es exponencial, en acuerdo con Floquet.
\end{ejemplo}

\begin{figure}[ht]
\centering
\includegraphics[width=0.85\linewidth]{cap09/f_yoshizawa}
\caption{La cadena del \cref{ej:vt-barbalat} desde $(1,0)$: $W=x_2^2\to0$ (LaSalle--Yoshizawa), luego $\dot x_2\to0$ (Barbalat) y, de la ecuación, $x_1=-(\dot x_2+x_2)/(2+\sin t)\to0$.}
\label{fig:vt-yoshizawa}
\end{figure}

\codigorepo{cap09_variantes/C9_Ej9_MathieuBarbalat.m}{yoshizawa}{Los cinco pasos, señal por señal}

\begin{receta}[Cuando $\Vdot\le-W(x)$ con $W$ solo semidefinida]
\begin{enumerate}
  \item Verifique que $V$ es definida positiva y decreciente. Entonces, por el teorema de LaSalle--Yoshizawa, las soluciones están acotadas y $W(x(t))\to0$.
  \item Identifique la variable que falta, $x_1$ en Mathieu o $\theta$ en el adaptativo, y busque una ecuación donde aparezca multiplicada por algo que no se anula.
  \item Derive una vez la variable que ya converge, verifique que la derivada segunda está acotada y aplique Barbalat para concluir que su derivada tiende a cero.
  \item Lea la variable que falta de la ecuación. Si el coeficiente puede anularse o desaparecer con el tiempo, el paso falla, y casi siempre la conclusión también.
\end{enumerate}
\end{receta}

El último punto es el de la excitación persistente. En el control adaptativo $\dot e=-e+\theta\omega(t)$, $\dot\theta=-e\omega(t)$, con $V=e^2+\theta^2$ se tiene $\Vdot=-2e^2$ y LaSalle--Yoshizawa da $e\to0$; el paso~4 da $\theta\,\omega(t)\to0$, que implica $\theta\to0$ solo si $\omega$ no se apaga. Con $\omega=\sin t$, $\theta\to0$; con $\omega$ que decae exponencialmente, $\theta$ se congela en un valor distinto de cero (\cref{ej:av-adaptativo}).

\begin{trampa}[label=tr:vt-lasalle]{``LaSalle con $\{\Vdot=0\}$ funciona igual para $\dot x=f(t,x)$''}
\plausible{El argumento parece el mismo: $V$ baja, las soluciones están acotadas, y se acercan al conjunto donde $\Vdot=0$.}
Sin invarianza de los conjuntos límite, el paso ``el mayor conjunto invariante dentro de $E$'' pierde su justificación. El \cref{ej:vt-e1} lo muestra: con $\dot x=-x/(1+t)^2$ y $V=x^2$, $\Vdot=-2x^2/(1+t)^2$ se anula solo en $x=0$, para todo $t$. Un ``LaSalle variante'' ingenuo daría $E=M=\{0\}$ y concluiría $x(t)\to0$; pero $x(t)\to x_0e^{-1/(1+t_0)}\neq0$.
\correcto{Para sistemas variantes en general use LaSalle--Yoshizawa y Barbalat; para sistemas periódicos en $t$, la versión periódica de LaSalle sí es válida, porque el mapa estroboscópico es invariante en el tiempo \citep[cap.~8]{RHL77}; \citep[\S8.3, nota~14]{Kha02}.}
\end{trampa}

\begin{profundiza}[El teorema de Matrosov]
Cuando una sola $V$ no da UAS, Matrosov~\citep{Mat62} combina una función $V$ con $\Vdot\le-W(x)\le0$ y una función auxiliar $U(t,x)$ acotada cuya derivada es estrictamente negativa cerca del conjunto $\{W=0\}$, fuera del origen. Las dos juntas dan estabilidad asintótica \emph{uniforme} sin exigir periodicidad. Para la Mathieu amortiguada, $U=x_1x_2$ cumple $\dot U=x_2^2-x_1x_2-(2+\sin t)x_1^2$, que en $x_2=0$ vale $-(2+\sin t)x_1^2<0$. Es la herramienta natural para sistemas mecánicos con amortiguamiento parcial y para el análisis de controladores adaptativos.
\end{profundiza}

%% ---------------------------------------------------------------------
\section{Inestabilidad}
\label{sec:vt-inestabilidad}
%% ---------------------------------------------------------------------

Hasta aquí todos los teoremas probaban estabilidad. Probar \emph{inestabilidad} también requiere una función, pero el argumento es el contrario: una cantidad que crece cerca del origen y que, por estar acotada en una bola, obliga a la solución a salir de ella.

\begin{definicion}[label=def:vt-inestable]{Inestabilidad}
El origen de \eqref{eq:vt-sistema} es \clave{inestable} si no es estable: existen $\varepsilon>0$ y $t_0\ge0$ tales que, para todo $\delta>0$, hay un $x(t_0)$ con $\norm{x(t_0)}<\delta$ y un $t\ge t_0$ con $\norm{x(t)}\ge\varepsilon$.
\end{definicion}

La negación es literal: basta \emph{un} $\varepsilon$, y para cada $\delta$ basta \emph{una} condición inicial que salga de $\Bola\varepsilon$. No se pide que todas las soluciones se alejen, ni que se vayan a infinito.

\begin{trampa}[label=tr:vt-inestable]{``Inestable quiere decir que las soluciones se van a infinito''}
\plausible{Es la imagen de $\dot x=x$, y la de los ejemplos de la \cref{sec:vt-ejemplos}.}
Para van der Pol, $\dot x_1=x_2$, $\dot x_2=-x_1+(1-x_1^2)x_2$, el origen es un foco inestable (valores propios $\tfrac12\pm j\tfrac{\sqrt3}2$), y \emph{todas} las soluciones están acotadas: las que salen de cerca del origen terminan en el ciclo límite, que cumple $1.53\le\norm x\le2.83$. Con $\varepsilon=1$ basta: desde $(\delta,0)$, la solución sale de $\Bola1$ en $t\approx2\ln(1/\delta)$ ($8.90$, $18.97$ y $27.23$ para $\delta=10^{-2},10^{-4},10^{-6}$), por pequeño que sea $\delta$ (\cref{fig:vt-vdp}).
\correcto{Inestable es la negación de estable: existe un $\varepsilon$ que, para todo $\delta$, alguna solución supera. Ni ``no asintóticamente estable'' implica inestable (un centro es estable), ni inestable implica no acotado.}
\end{trampa}

\begin{figure}[ht]
\centering
\includegraphics[width=\linewidth]{cap09/f_vdp}
\caption{Van der Pol con $\mu=1$. Izquierda: soluciones desde cerca del origen y desde fuera del ciclo; todas acotadas, todas al ciclo. La circunferencia punteada es $\norm x=1$. Derecha: $\norm{x(t)}$ desde $(\delta,0)$ en escala logarítmica; el tiempo para salir de $\Bola1$ crece como $2\ln(1/\delta)$, pero siempre es finito.}
\label{fig:vt-vdp}
\end{figure}

\codigorepo{cap09_variantes/C9_Ej11_VanDerPol.m}{vdp}{Inestable y acotado: tiempos de salida de $\Bola1$}

\subsection{El teorema de Chetaev}

Para sistemas invariantes en el tiempo, el resultado más útil es el de Chetaev \citep{Che61}, \citep[Teo.~4.3]{Kha02}.

\begin{teorema}[label=thm:vt-cetaev]{Chetaev}
Sea $\dot x=f(x)$ con $f(0)=0$ y $V:\D\to\R$ de clase $\Cont1$ con $V(0)=0$ y tal que $V(x_0)>0$ para algún $x_0$ de norma arbitrariamente pequeña. Sea $r>0$ con $\overline{\Bola r}\subset\D$ y
\[
  U=\{x\in\overline{\Bola r}:\ V(x)>0\} .
\]
Si $\Vdot(x)>0$ en $U$, el origen es inestable.
\end{teorema}

\begin{demostracion}
Sea $x_0\in U$ con $\norm{x_0}$ tan pequeño como se quiera, y $a=V(x_0)>0$. Mientras la solución esté en $U$, $V$ crece, así que $V(x(t))\ge a$. La solución no puede salir de $U$ por la frontera $\{V=0\}$, porque allí $V$ valdría $0<a$; solo puede salir de la bola. Considere el compacto $\Gamma=\{x\in\overline{\Bola r}:\ V(x)\ge a\}$, contenido en $U$ (no contiene al origen, donde $V=0$), y $\gamma=\min_\Gamma\Vdot>0$. Mientras $x(t)\in\Gamma$, $V(x(t))\ge a+\gamma t$, que crece sin cota; pero $V$ está acotada en $\overline{\Bola r}$. Luego la solución sale de $\Bola r$, y como $x_0$ era arbitrariamente cercano al origen, este es inestable con $\varepsilon=r$.
\end{demostracion}

$V$ \emph{no} tiene que ser definida positiva: basta con que sea positiva en algún sector que llega hasta el origen, y el teorema solo mira lo que pasa en ese sector. En la silla $\dot x_1=x_1$, $\dot x_2=-x_2$, $V=\tfrac12(x_1^2-x_2^2)$ es indefinida y $\Vdot=x_1^2+x_2^2>0$ en todas partes.

\begin{ejemplo}[label=ej:vt-cetaev]{Un valor propio cero}
Para $\dot x_1=x_1^2+x_2^3$, $\dot x_2=-x_2+x_1^3$, pruebe que el origen es inestable.
\solucion La linealización tiene valores propios $0$ y $-1$: el método indirecto no decide. Tome $V=x_1-\tfrac12x_2^2$, de modo que $U=\{x_1>x_2^2/2\}\cap\Bola r$ es una región con forma de parábola que llega al origen (en $(\delta,0)$, $V=\delta>0$). Entonces
\[
  \Vdot=(x_1^2+x_2^3)-x_2(-x_2+x_1^3)=x_1^2+x_2^2+x_2^3-x_1^3x_2 .
\]
En $\Bola r$, $\abs{x_2^3}\le r\,x_2^2$ y $\abs{x_1^3x_2}\le x_1^2\,\abs{x_1x_2}\le\tfrac{r^2}2x_1^2$, así que
\[
  \Vdot\ge\Bigl(1-\frac{r^2}2\Bigr)x_1^2+(1-r)\,x_2^2\ge\min\Bigl\{1-\frac{r^2}2,\ 1-r\Bigr\}\norm x^2 .
\]
Con $r=\tfrac12$, $\Vdot\ge\tfrac12\norm x^2>0$ en $\Bola{1/2}\setminus\{0\}$, en particular en $U$: el origen es inestable. La constante $\tfrac12$ es exacta: el mínimo de $\Vdot/\norm x^2$ en $\Bola{1/2}$, calculado numéricamente, es $0.50$, y se alcanza cerca de $(0,-\tfrac12)$.
\end{ejemplo}

\begin{trampa}[label=tr:vt-cetaevx1]{``Con $V=x_1$ basta: $\dot x_1=x_1^2+x_2^3$ es casi siempre positivo''}
\plausible{$V=x_1$ es más simple, $U=\{x_1>0\}$ llega al origen y $x_1^2>0$ domina cerca del eje $x_1$.}
En $x=(\delta^2,-\delta)$, con $0<\delta<1$, $V=\delta^2>0$ y $\Vdot=\delta^4-\delta^3<0$. Esos puntos están en $U$ arbitrariamente cerca del origen, así que $\Vdot$ no es positiva en $U$ para ningún $r$.
\correcto{Una $V$ de Chetaev se construye para que la parte buena domine a la mala. El término $-\tfrac12x_2^2$ aporta $-x_2\dot x_2=x_2^2-x_1^3x_2$, y ese $x_2^2$ domina a $x_2^3$; además, recorta $U$ a la parábola $x_1>x_2^2/2$, donde $x_1$ no es pequeño frente a $x_2$.}
\end{trampa}

\subsection{Dos teoremas de inestabilidad para $\dot x=f(t,x)$}

Para sistemas variantes en el tiempo hay dos versiones, que se prueban con el mismo argumento de Chetaev \citep{SlLi91}. En las dos, $V:[0,\infty)\times\Bola r\to\R$ es de clase $\Cont1$, $V(t,0)=0$, es decreciente ($\abs{V(t,x)}\le W_2(x)$) y, para algún $t_0$, $V(t_0,\cdot)$ toma valores positivos en puntos arbitrariamente cercanos al origen.

\begin{teorema}[label=thm:vt-inestable1]{Primer teorema de inestabilidad}
Si además $\Vdot(t,x)\ge W(x)$ en $[0,\infty)\times\Bola r$, con $W$ continua y definida positiva, el origen es inestable.
\end{teorema}

\begin{teorema}[label=thm:vt-inestable2]{Segundo teorema de inestabilidad}
Si además existe $\lambda>0$ con
\[
  \Vdot(t,x)-\lambda\,V(t,x)\ge0\qquad\text{en }[0,\infty)\times\Bola r ,
\]
el origen es inestable.
\end{teorema}

\begin{demostracion}
Sea $x_0$, de norma tan pequeña como se quiera, con $a=V(t_0,x_0)>0$, y suponga que la solución se queda en $\Bola r$. En el segundo teorema, $\frac{\dif}{\dif t}\bigl(e^{-\lambda(t-t_0)}V\bigr)=e^{-\lambda(t-t_0)}(\Vdot-\lambda V)\ge0$, así que $V(t,x(t))\ge ae^{\lambda(t-t_0)}\to\infty$. En el primero, $V$ crece, $V\ge a$, luego $W_2(x(t))\ge a$ y $\norm{x(t)}\ge\gamma>0$ para algún $\gamma$; en el anillo $\gamma\le\norm x\le r$, $W\ge m>0$ y $V\ge a+m(t-t_0)\to\infty$. En los dos casos se contradice $V\le\max_{\overline{\Bola r}}W_2$.
\end{demostracion}

\begin{ejemplo}[label=ej:vt-inestable1]{Primer teorema donde la linealización calla}
Para $\dot x_1=2x_2+x_1(x_1^2+2x_2^4)$, $\dot x_2=-2x_1+x_2(x_1^2+x_2^4)$, con $V=\tfrac12\norm x^2$,
\[
  \Vdot=x_1^4+x_1^2x_2^2+2x_1^2x_2^4+x_2^6 ,
\]
que es definida positiva (todos los términos son no negativos, y $x_1^4+x_2^6>0$ si $x\neq0$). El origen es inestable. La linealización, $\dot x_1=2x_2$, $\dot x_2=-2x_1$, es un centro con valores propios $\pm2j$, y no dice nada.
\end{ejemplo}

\begin{ejemplo}[label=ej:vt-inestable2]{Segundo teorema: $\Vdot$ no tiene que ser positiva}
Para
\[
  \dot x_1=x_1+3x_2\sin^2x_2+5x_1x_2^2\sin^2x_1,\qquad
  \dot x_2=3x_1\sin^2x_2+x_2-5x_1^2x_2\cos^2x_1 ,
\]
con $V=\tfrac12(x_1^2-x_2^2)$, los términos con $3\sin^2x_2$ se cancelan y
\[
  \Vdot=x_1^2-x_2^2+5x_1^2x_2^2\bigl(\sin^2x_1+\cos^2x_1\bigr)=2V+5x_1^2x_2^2 .
\]
Entonces $\Vdot-2V=5x_1^2x_2^2\ge0$, y con $\lambda=2$ el origen es inestable. Note que $\Vdot$ \emph{no} es positiva: en $(0,0.5)$, $V=-0.125$ y $\Vdot=-0.25$. El primer teorema no aplica; el segundo sí, porque solo pide que $V$ crezca exponencialmente donde ya es positiva.
\end{ejemplo}

\begin{figure}[ht]
\centering
\includegraphics[width=\linewidth]{cap09/f_inestabilidad}
\caption{Izquierda: \cref{ej:vt-inestable1}; la solución desde $(0.25,0)$ gira y se aleja. Centro: \cref{ej:vt-inestable2}; en coral la región $V>0$ ($\abs{x_1}>\abs{x_2}$), de donde las soluciones salen. Derecha: \cref{ej:vt-cetaev}; en verde $U=\{x_1>x_2^2/2\}$ dentro de $\Bola{1/2}$ (punteada); las soluciones que nacen en $U$ salen de la bola.}
\label{fig:vt-inestabilidad}
\end{figure}

\codigorepo{cap09_variantes/C9_Ej10_Inestabilidad.m}{inestabilidad}{Las tres $\Vdot$ y la cota de Chetaev en una malla}

\video{V9_4_cetaev}{La región de Chetaev: donde $V>0$, $V$ solo puede crecer}

%% ---------------------------------------------------------------------
\section{Para llevar}
\label{sec:vt-cierre}
%% ---------------------------------------------------------------------

\begin{parallevar}
\begin{enumerate}
  \item \textbf{La hora de arranque importa:} para $\dot x=f(t,x)$ la solución depende de $t_0$, y aparecen las nociones uniformes. US: $\delta$ no depende de $t_0$. UAS: tampoco el tiempo de convergencia. En sistemas invariantes en el tiempo (y periódicos) la uniformidad es gratis.
  \item \textbf{Funciones de comparación:} UAS es lo mismo que una sola cota $\norm{x(t)}\le\beta(\norm{x(t_0)},t-t_0)$ con $\beta\in\KL$. Hay que revisar las dos variables de $\beta$: estrictamente creciente en el tamaño, decreciente a cero en el tiempo transcurrido.
  \item \textbf{Lyapunov con $V(t,x)$:} $W_1(x)\le V\le W_2(x)$ y $\Vdot\le-W_3(x)$ dan UAS; con $W_1$ radialmente no acotada, GUAS; con cotas $k_i\norm x^a$, exponencial. ``Decreciente'' ($V\le W_2$) es lo que hace uniforme el resultado, y $\Vdot$ incluye $\partial V/\partial t$.
  \item \textbf{$\Vdot<0$ no basta:} hace falta una cota $-W_3(x)$ independiente de $t$. Tampoco bastan los valores propios congelados (Markus--Yamabe) ni una linealización sin uniformidad.
  \item \textbf{Con $\Vdot$ semidefinida:} LaSalle--Yoshizawa da $W(x(t))\to0$; Barbalat sobre la ecuación de la variable que falta completa la convergencia; en el caso lineal periódico, Floquet decide.
  \item \textbf{Inestabilidad:} es la negación de estabilidad, no ``escape a infinito''. Chetaev (invariante) y los dos teoremas de inestabilidad (variantes) piden una $V$ que sea positiva cerca del origen y crezca donde es positiva; $V$ puede ser indefinida.
\end{enumerate}
\end{parallevar}

\begin{trampa}[title={\faBug\hspace{0.45em}Las trampas del capítulo, en una caja}]{}
\begin{itemize}
  \item Converger no es converger uniformemente (\cref{tr:vt-converge}); un $\delta(\varepsilon,t_0)$ no es un $\delta(\varepsilon)$ (\cref{tr:vt-estable}).
  \item Decaer a cero en $s$ no basta para ser de clase $\KL$ (\cref{tr:vt-kl}).
  \item $\Vdot$ lleva el término $\partial V/\partial t$ (\cref{tr:vt-parcial}); $\Vdot<0$ en cada instante no es $\Vdot\le-W_3(x)$ (\cref{tr:vt-negativa}).
  \item Los valores propios de $A(t)$ no deciden (\cref{tr:vt-congelar}); la linealización necesita un jacobiano uniforme en $t$ (\cref{ej:vt-linealizacion}).
  \item LaSalle no se traslada tal cual a sistemas variantes (\cref{tr:vt-lasalle}).
  \item Inestable no es no acotado (\cref{tr:vt-inestable}); una $V$ de Chetaev se construye, no se adivina (\cref{tr:vt-cetaevx1}).
  \item Un certificado que falla no prueba inestabilidad: con la $g$ rápida del \cref{ej:vt-ganancia} el sistema sigue siendo exponencialmente estable (\cref{prob:vt-grapida}).
\end{itemize}
\end{trampa}

%% ---------------------------------------------------------------------
\section{Problemas}
\label{sec:vt-problemas}
%% ---------------------------------------------------------------------

Los problemas marcados con \etiquetaM{} requieren \matlab/\simulink; los marcados con \etiquetaEstrella{} son de profundización. Cada uno tiene una trampa: búsquela antes de dar la respuesta.

\subsection*{Uniformidad y funciones de comparación}

\begin{problema}[una tasa que se apaga más despacio]\label{prob:vt-raiz}
Para $\dot x=-x/\sqrt{1+t}$: (a) resuelva la ecuación; (b) calcule el tiempo de asentamiento $T(0.1,t_0)$ y evalúelo en $t_0=0$ y $t_0=100$; (c) ¿es el origen UAS?
\end{problema}

\begin{problema}[estable, pero no uniformemente]\label{prob:vt-6tsin}
\etiquetaM{} Para $\dot x=(6t\sin t-2t)\,x$: (a) escriba la solución; (b) con $t_0=2n\pi$, muestre que $x(t_0+\pi)/x(t_0)=e^{(4n+1)(6-\pi)\pi}$; (c) concluya que el origen no es uniformemente estable. (d) Simule desde $x(t_0)=10^{-3}$ con $n=0,1,2$ y grafique en escala logarítmica.
\end{problema}

\begin{problema}[\texorpdfstring{$\KL$}{KL} y flujos]\label{prob:vt-flujo}
Una función $\sigma\in\KL$ generada por $\dot y=-\alpha(y)$ cumple $\sigma(r,0)=r$ y $\sigma(\sigma(r,s_1),s_2)=\sigma(r,s_1+s_2)$. (a) De las funciones de la pausa de la \cref{sec:vt-comparacion}, ¿cuáles son generadas por una ecuación de comparación? (b) Muestre que $re^{-s}$ y $r/(rs+1)$ sí lo son, y recupere $\alpha$ con $\alpha(r)=-\partial\sigma/\partial s\,(r,0)$.
\end{problema}

\subsection*{Teoremas de Lyapunov}

\begin{problema}[decreciente o no {\citep[ejercicio~2.7]{KeBr23}}]\label{prob:vt-decreciente}
Para $V_1=x_1^2(1+\sin^2t)+x_2^2(1+\cos^2t)$ y $V_2=x_1^2+x_2^2(1+t)$, escriba $\alpha_1$ y, si existe, $\alpha_2$. Explique en una frase por qué no existe para $V_2$.
\end{problema}

\begin{problema}[atracción sin exponencial]\label{prob:vt-cubico}
Para $\dot x=-(1+g(t))\,x^3$ con $0\le g(t)\le k$: (a) con $V=x^2/2$ pruebe que el origen es GUAS; (b) pruebe que no es exponencialmente estable. ¿Para qué sirve la cota $g\le k$?
\end{problema}

\begin{problema}[cota cruzada {\citep[ejercicio~4.2]{Ser00}}]\label{prob:vt-seron2}
Para $\dot x_1=-2x_1+g(t)x_2$, $\dot x_2=g(t)x_1-2x_2$ con $\abs{g(t)}\le1$, pruebe que el origen es GES y dé la tasa. ¿Hace falta alguna condición sobre $\dot g$?
\end{problema}

\begin{problema}[rotación y disipación cúbica {\citep[ejercicio~4.3]{Ser00}}]\label{prob:vt-seron3}
Para $\dot x_1=x_2-g(t)x_1\norm x^2$, $\dot x_2=-x_1-g(t)x_2\norm x^2$ con $k\le g(t)\le M$, $k>0$: (a) pruebe que el origen es GUAS; (b) pruebe que no es exponencialmente estable; (c) ¿qué dice la linealización?
\end{problema}

\begin{problema}[la ecuación de Lyapunov diferencial, escalar]\label{prob:vt-lyapdif}
\etiquetaEstrella{} Para $\dot x=-a(t)x$ con $a(t)=2+\sin t$: (a) escriba la ecuación \eqref{eq:vt-lyapdif} con $Q=1$; (b) muestre que $p(t)=\int_t^\infty e^{-2\int_t^\tau a(\sigma)\dif\sigma}\dif\tau$ la resuelve; (c) pruebe que $\tfrac16\le p(t)\le\tfrac12$. \etiquetaM{} (d) Calcule numéricamente el rango de $p$.
\end{problema}

\subsection*{Sistemas lineales variantes y ejemplos}

\begin{problema}[Markus--Yamabe a mano]\label{prob:vt-markus}
Verifique la segunda columna de $\Phi(t,0)$ en el \cref{ej:vt-markus}, compruebe la fórmula de Liouville $\det\Phi(t,0)=\exp\int_0^t\operatorname{tr}A$, y encuentre una condición inicial cuya solución tienda a cero.
\end{problema}

\begin{problema}[la ganancia rápida]\label{prob:vt-grapida}
\etiquetaM{} En el \cref{ej:vt-ganancia} con $g=(1+\sin10t)/2$: (a) verifique que $\dot g\le g$ falla; (b) calcule numéricamente la matriz de monodromía con periodo $\pi/5$ y sus multiplicadores; (c) ¿es estable el origen? ¿Qué dice eso de la $V$ del ejemplo? (d) Explique el módulo de los multiplicadores con la fórmula de Liouville.
\end{problema}

\subsection*{Inestabilidad}

\begin{problema}[¿hasta dónde llega la cota de Chetaev?]\label{prob:vt-cetaevr}
En el \cref{ej:vt-cetaev}, ¿para qué valores de $r$ la cota $\min\{1-r^2/2,\ 1-r\}\norm x^2$ prueba que $\Vdot>0$ en $\Bola r\setminus\{0\}$? ¿Cambia eso la conclusión de inestabilidad?
\end{problema}

\begin{problema}[van der Pol: qué teorema sirve]\label{prob:vt-vdp}
\etiquetaEstrella{} Para van der Pol con $\mu=1$: (a) con $V=\tfrac12\norm x^2$, calcule $\Vdot$ y explique por qué el primer teorema de inestabilidad no aplica directamente; (b) pruebe la inestabilidad con el método indirecto; (c) explique por qué el tiempo de salida de $\Bola1$ desde $(\delta,0)$ es aproximadamente $2\ln(1/\delta)$.
\end{problema}

%% ---------------------------------------------------------------------
\section{Soluciones seleccionadas}
\label{sec:vt-soluciones}
%% ---------------------------------------------------------------------

\paragraph{Problema~\ref{prob:vt-raiz}}
(a) $x(t)=x_0\exp\bigl(-2(\sqrt{1+t}-\sqrt{1+t_0})\bigr)$. (b) $x=0.1x_0$ cuando $\sqrt{1+t}=\sqrt{1+t_0}+\tfrac12\ln10$, es decir $T=\ln10\,\sqrt{1+t_0}+\tfrac14\ln^210$: $T=3.63$ con $t_0=0$ y $T=24.47$ con $t_0=100$. (c) No: $T$ crece como $\sqrt{t_0}$. Es US ($\abs{x(t)}\le\abs{x_0}$) y AS, pero no UAS.

\paragraph{Problema~\ref{prob:vt-6tsin}}
(a) $x(t)=x(t_0)\exp\bigl(6\sin t-6t\cos t-t^2-6\sin t_0+6t_0\cos t_0+t_0^2\bigr)$. (b) Entre $2n\pi$ y $(2n+1)\pi$, $\int6\tau\sin\tau=6(4n+1)\pi$ y $\int2\tau=(4n+1)\pi^2$. (c) Para tener $\abs{x(t)}<\varepsilon$ desde $t_0=2n\pi$ hace falta $\delta\le\varepsilon e^{-(4n+1)(6-\pi)\pi}$, que tiende a cero con $n$: no hay $\delta(\varepsilon)$ uniforme. El exponente total tiende a $-\infty$, así que el origen es AS.

\paragraph{Problema~\ref{prob:vt-flujo}}
(a) Ninguna: (a) y (d) son $\KL$ pero $\sigma(r,0)=r^2$ y $2r$, distintos de $r$; (b) y (c) no son $\KL$. (b) $re^{-s}$ viene de $\dot y=-y$ y $r/(rs+1)$ de $\dot y=-y^2$; $-\partial_s(re^{-s})|_{s=0}=r$ y $-\partial_s\bigl(r/(rs+1)\bigr)|_{s=0}=r^2$.

\paragraph{Problema~\ref{prob:vt-decreciente}}
$V_1$: $\alpha_1(r)=r^2$ y $\alpha_2(r)=2r^2$, porque $1\le1+\sin^2t\le2$. $V_2$: $\alpha_1(r)=r^2$; en $x=(0,1)$, $V_2=1+t$, que no está acotado por ningún $\alpha_2(1)$.

\paragraph{Problema~\ref{prob:vt-cubico}}
(a) $\Vdot=-(1+g)x^4\le-x^4$, con $V$ invariante en el tiempo, definida positiva y radialmente no acotada: GUAS, con $\abs{x(t)}\le\abs{x_0}/\sqrt{1+2x_0^2(t-t_0)}$. (b) Como $\dot x\ge-(1+k)x^3$ para $x>0$, por comparación $x(t)\ge x_0/\sqrt{1+2(1+k)x_0^2(t-t_0)}$, que decae como $1/\sqrt t$. La cota $g\le k$ no se usa para GUAS: solo para negar la estabilidad exponencial.

\paragraph{Problema~\ref{prob:vt-seron2}}
Con $V=\norm x^2$, $\Vdot=-4\norm x^2+4g\,x_1x_2=x\T\bigl[\begin{smallmatrix}-4&2g\\2g&-4\end{smallmatrix}\bigr]x\le-2\norm x^2$, porque los valores propios son $-4\pm2g\le-2$. Con $k_1=k_2=1$, $k_3=2$, $a=2$: $\norm{x(t)}\le\norm{x(t_0)}e^{-(t-t_0)}$. No hace falta nada sobre $\dot g$: $V$ no depende de $t$.

\paragraph{Problema~\ref{prob:vt-seron3}}
(a) Con $V=\norm x^2$ los términos de rotación se cancelan y $\Vdot=-2g\norm x^4\le-2k\norm x^4$: GUAS. (b) Con $g\le M$, $\frac{\dif}{\dif t}\norm x^2\ge-2M\norm x^4$ y $\norm{x(t)}^2\ge\norm{x_0}^2/(1+2M\norm{x_0}^2(t-t_0))$, que decae como $1/t$. (c) Un centro ($\pm j$): no decide.

\paragraph{Problema~\ref{prob:vt-lyapdif}}
(a) $\dot p-2a(t)p=-1$. (b) Derivando la integral respecto a su límite inferior y a $t$ dentro del exponente, $\dot p=-1+2a(t)p$. (c) Como $1\le a\le3$, $e^{-6(\tau-t)}\le e^{-2\int_t^\tau a}\le e^{-2(\tau-t)}$, e integrando, $\tfrac16\le p\le\tfrac12$. (d) $0.167\le p(t)\le0.467$.

\paragraph{Problema~\ref{prob:vt-markus}}
Para $\psi=e^{-t}(s,c)$: $\dot\psi=e^{-t}(c-s,\ -s-c)$ y $A\psi$ da lo mismo. $\det\Phi=e^{t/2}e^{-t}(c^2+s^2)=e^{-t/2}=e^{\int_0^t(-1/2)}$. Desde $x(0)=(0,1)$, $x(t)=e^{-t}(\sin t,\cos t)\to0$: la inestabilidad no es de todas las soluciones, sino de las que tienen componente en la primera columna.

\paragraph{Problema~\ref{prob:vt-grapida}}
(a) En $t=0$, $g=\tfrac12$ y $\dot g=5$. (b) Los multiplicadores son $0.482\pm0.230j$, con módulo $0.533=e^{-\pi/5}$. (c) Sí, globalmente exponencialmente estable, con tasa $1$: la $V$ del ejemplo dejó de ser un certificado, pero el sistema no cambió de naturaleza. (d) $\operatorname{tr}A=-2$ para todo $t$, así que $\det M=e^{-2\pi/5}$, y con multiplicadores complejos conjugados, $\abs\mu=e^{-\pi/5}$.

\paragraph{Problema~\ref{prob:vt-cetaevr}}
Para $r<1$, porque $1-r>0$ y $1-r^2/2>0$. No: la inestabilidad es local y basta un $r$.

\paragraph{Problema~\ref{prob:vt-vdp}}
(a) $\Vdot=(1-x_1^2)x_2^2\ge0$ cerca del origen, pero se anula en todo el eje $x_2=0$: no es definida positiva. (b) El jacobiano en el origen tiene valores propios $\tfrac12\pm j\tfrac{\sqrt3}2$, con parte real positiva. (c) Cerca del origen la solución es aproximadamente la de la linealización, $\norm{x(t)}\approx\delta e^{t/2}$, que vale $1$ en $t=2\ln(1/\delta)$; el término $O(1)$ corrige la parte en que la no linealidad ya pesa.

\respuestaspausas

\begin{parallevar}[Nota de transparencia]
En la redacción de este capítulo y en la escritura y corrección de los códigos que lo acompañan se contó con la asistencia de Claude, un modelo de inteligencia artificial desarrollado por Anthropic. El contenido fue revisado y validado por el autor, quien asume la responsabilidad por lo aquí expuesto.
\end{parallevar}

